% concatenated from high-d_JAT.tex
\documentclass{article}

%%%%%%%%%%%%
%% THEMES %%
%%%%%%%%%%%%


\usepackage{a4}
\usepackage{graphicx}
\usepackage{subfigure}
\RequirePackage{amsthm,amsmath,amssymb,amsfonts,enumerate}

\usepackage{tikz}
\usepackage{makeidx}         % allows index generation
\usepackage{graphicx}        % standard LaTeX graphics tool
                             % when including figure files
\usepackage{multicol}        % used for the two-column index

\usepackage{amsmath}
\usepackage{amsfonts}
\usepackage{amsbsy}
\usepackage{graphics}
\usepackage{epsfig}
\usepackage{amssymb}
\usepackage{mathrsfs}
\usepackage{dsfont}

\usepackage{url}

\usepackage{color}
\usepackage[dvipsnames, table]{xcolor}
%\usepackage[textsize=scriptsize]{todonotes}
%
%%\usepackage{showlabels}
%
\usepackage{comment}


\def\Px{\mathsf{P}}
\def\mgbb{{\boldsymbol{\gamma}}}

\def\eqd{\stackrel{{d}}{=}}

\def\deltau{\underline{\delta}}

\newcommand{\HD}{\textcolor{red}}

\newcommand{\AZ}{\textcolor{blue}}

\newcommand{\W}{\textcolor{white}}

\newcommand{\Luccomment}[1]{\todo[color=yellow]{#1}}
\newcommand{\Lucinline}[1]{\todo[inline,color=yellow]{#1}}
\newcommand{\New}{\color{ForestGreen} }
\newcommand{\Luc}{\color{magenta} }
\newcommand{\Chk}{\color{red} }

%_______________________________________________________________
%%%%%%%%%%%%%%%%%%%%%%%%%%%%%%%%%%%%%%%%%%%%%%%%%%%%%%
% our Latex commands and shortcuts
%
%\input{Latex-commands}
\def\RR{\mathds{R}}
\def\QQ{\mathds{Q}}
\def\NN{\mathds{N}}
\def\SX{{\mathscr X}}
\def\Ib{\mathbf{I}}
\def\xb{\mathbf{x}}
\def\Xb{\mathbf{X}}
\def\Ex{\mathsf{E}}
\def\var{\mathsf{var}}
\def\simd{\stackrel{d}{\sim}}
\def\me{\epsilon}
\def\ml{\lambda}
\def\ms{\sigma}
\newcommand{\vsp} {\vspace{0.3cm}}
\def\PP{\mathbb{P}}
\def\SB{{\mathscr B}}
\def\SS{{\mathcal S}}
\def\0b{\mathbf{0}}
\newcommand{\bea}{\begin{eqnarray*}}
\newcommand{\eea}{\end{eqnarray*}}
\newcommand{\be}{\begin{eqnarray}}
\newcommand{\ee}{\end{eqnarray}}
\def\Rb{\mathbf{R}}
\def\ub{\mathbf{u}}
\def\CR{\mathsf{CR}}
\def\mg{\gamma}
\def\dd{\mbox{\rm d}}
\def\rad{\stackrel{\rm d}{\ra}}
\def\ra{\rightarrow}
\def\SN{{\mathcal N}}
\def\ma{\alpha}
\def\e1{\mathrm{e}}
\newcommand{\fin} {\mbox{}~\hfill{\lower-0.3ex\hbox{$\triangleleft$}}}
\def\Ind{\mathds{1}}
\def\SO{{\mathcal O}}
\def\SH{{\mathcal H}}
\def\SV{\mathscr{V}}
\def\bb{\mathbf{b}}



\newtheorem{theorem}{Theorem}[section]
\newtheorem{proposition}{Proposition}[section]
\newtheorem{lemma}{Lemma}[section]
\newtheorem{corollary}{Corollary}[section]
\newtheorem{definition}{Definition}[section]
\newtheorem{remark}{Remark}[section]
\newtheorem{example}{Example}
%


\begin{document}

%\title{quantisation in high dimensional sets}

\title{Non-asymptotic quantisation of spherically symmetric distributions}
%\title{Towards optimal quantisation of spherically symmetric distributions in high dimensions: a non-asymptotic approach}
%\title{Towards optimal non-asymptotic quantisation\\ of spherically symmetric distributions\\ in high dimensions}

%\title{Incremental space-filling designs in high dimension based on quantisation}

%\author{Luc Pronzato$^{1}$ and Anatoly Zhigljavsky$^{2}$}
%\address{ $^{1}$Universit\'e C\^ote d'Azur-CNRS, Sophia Antipolis, France, {\tt pronzato@i3s.unice.fr}\\
%\mbox{}\\
%$^{2}$School of Mathematics, Cardiff University, UK, {\tt ZhigljavskyAA@cardiff.ac.uk} }

\author{Luc Pronzato\footnote{(corresponding author) Laboratoire I3S, CNRS-Universit\'e C\^ote d'Azur, Sophia Antipolis, France, {\tt pronzato@i3s.unice.fr}} \ and Anatoly Zhigljavsky\footnote{School of Mathematics, Cardiff University, UK, {\tt ZhigljavskyAA@cardiff.ac.uk}}}

\date{\today}

\maketitle

\begin{abstract}
Zador’s celebrated theorem is a cornerstone of optimal quantisation, establishing both the weak limit of the empirical distribution of an $n$-point optimal quantiser in $\RR^d$ and the decay rate of the associated $L_s$-mean quantisation error. However, for large dimensions $d$, observing this asymptotic behaviour demands an astronomically large sample size $n$, which grows super-exponentially with $d$. Through a detailed analysis of the quantisation problem for spherically symmetric distributions, we demonstrate that for moderate $n$ random quantisers uniformly distributed on a sphere of suitable radius $R$ achieve exceptional performance. The expected distortion, expressed as a triple integral, can be computed with arbitrary precision, and the optimal radius $R$ can be efficiently determined numerically. Leveraging results from extreme-value theory, we derive approximations for $R$, particularly in scenarios where $n$ scales with $d$. Depending on the growth rate of $n$, $R$ may either converge to zero or approach a limiting value $R_\infty$ that is independent of $s$.
\end{abstract}

{\bf keywords:} quantisation, distortion, spherically symmetric distribution, Zador's theorem, extreme-value theory

%{\tableofcontents}

%\footnote{\AZ {the whole paper is about spherically symmetric distributions. I suggest to remove "spherically symmetric distribution" from section/subsection titles}}


%\footnote{\AZ {JAT is ok but the paper looks more suitable for something like  Journal of Multivariate Analysis or one of SIAM journals}}

%\footnote{\AZ  {Can we move some figures (plus related explanations) to the very end of paper? They kind of distract the main lines. The conclusions from the numerical results would be enough in the main text. } }

\section{Introduction}
\label{sec:intr}

For $\mu$ a measure on $\SX\subseteq\RR^d$ (or a $d$-dimensional Riemannian manifold),
with density $\varphi$ and finite moment of order $s'>s\geq 1$, the $(\mu,s)$-distortion on an $n$-point set $\Xb_n=\{\xb_1,\ldots,\xb_n\}\in\RR^{n\times d}$ is $D_{\mu,s}(\Xb_n)=\Ex_\mu\{ \min_{i=1,\ldots,n} \|U-\xb_i\|^s \}$, where $U\simd\mu$. Zador's celebrated theorem \cite{Zador82} states that the empirical distribution of an $n$-point optimal quantiser $\Xb_n^*$ minimising $D_{\mu,s}(\Xb_n)$ converges weakly to the probability measure with density proportional to $\varphi^{d/(d+s)}$ as $n\to\infty$, and the normalised optimal $(\mu,s)$-distortion $n^{1/d}\,D_{\mu,s}(\Xb_n^*)$ reaches asymptotically its minimum as $n\to\infty$. The normalisation by $n^{1/d}$ is essential, in the sense that these asymptotic considerations are not applicable if $n^{1/d}$  is not large enough. For large $d$, this requires astronomically large values of $n$ to approach the asymptotic behaviour.

The purpose of this paper is to show that for $\mu$ spherically symmetric, random quantisers with a suitable spherically symmetric distribution can achieve exceptional performance. When $n$ is large but $n^{1/d}$ is small, the best random quantisers we obtain are for designs $\Xb_n$ uniform on a sphere $\SS_{d-1}(R)$ with suitable radius $R$, and the variability of $D_{\mu,s}(\Xb_n)$ vanishes as $n$ increases. Using extreme-value theory, we derive approximations for the radius $R$ of $\SS_{d-1}$ and for the best expected $(\mu,s)$-distortion. In the particular case $s=2$, $R$ is proportional to $\Ex_\mu\{\|U\|\}$ with a constant depending only on $n$ and $d$, and the best expected $(\mu,2)$-distortion is $\Ex_\mu\{\|U\|^2\}-R^2$; see Corollary~\ref{Coro:astar-s2-Pa}. We then identify different asymptotic regimes where $d$ increases and $n$ grows with $d$. When $\mu$ possesses a norm-concentration property and concentrates on a sphere of radius $r$ as $d \to \infty$, the optimal radius $R$ of the random quantiser exhibits distinct behaviours depending on the growth rate of $n$ relative to $d$:
for a sub-exponential growth, $R \to 0$; for a super-exponential growth, $R \to r$; for an exponential growth $n = \ml^d(1 + o(1))$, $R \to r\sqrt{1 - 1/\ml^2}$; see Proposition~\ref{Prop:extreme-value-general-mu}.
Overall, the paper provides a framework for generating nearly optimal quantisers in the large $n$ and $d$ situation for spherically symmetric distributions.
Sequences of nested designs with low quantisation error can thus be efficiently generated, offering a compelling alternative to the classical greedy-packing algorithm for space-filling design (see, e.g., \cite{PZ2022}). While this paper focuses exclusively on spherically symmetric distributions, the conclusion highlights that extending this approach to the more common case of quantising the uniform measure in the $d$-dimensional hypercube yields promising numerical results.

\vsp
The paper is organised as follows. Section~\ref{S:distancecdt-etc} introduces the key concepts used throughout the paper. It also derives an explicit integral expression of the expected $(\mu,s)$-distortion of a random quantiser whose $n$ points are independently identically distributed (i.i.d.) according to a spherically symmetric distribution $\PP$. Section~\ref{S:ball-sphere-annulus} examines three emblematic cases: (\textit{i}) $\mu$ is uniform on the unit sphere, (\textit{ii}) $\mu$ is uniform in the unit ball, and (\textit{iii}) $\mu$ follows a multivariate spherically symmetric normal distribution. In the latter two scenarios, in addition to the case where $\PP$ is uniform on a sphere, we also analyse the cases where $\PP$ is uniform in a ball and where $\PP$ is a multivariate spherically symmetric normal distribution. Section~\ref{S:asymptotic} applies extreme-value theory to derive asymptotic results ($n \to \infty$) for random quantisers uniform on a sphere. The analysis is conducted in two scenarios: when $\mu$ is uniform on a sphere; when $\mu$ is spherically symmetric and satisfies a norm-concentration property, with the ball-uniform and normal distributions serving as key examples.
Section~\ref{S:conclusion} provides a brief conclusion to the paper.

\vsp
We denote by $\SB_d(r)$ the $d$-dimensional (closed) Euclidean ball with center $\0b_d$ and radius $r$ %, $\SB_d^o(r)$ is the interior of $\SB_d(r)$.
and by $\SS_{d-1}(r)$ the $d$-dimensional Euclidean sphere with center $\0b_d$ and radius $r$; $\|\cdot\|$ is the $\ell_2$-norm.


%___________________________________________________________________________
\section{Distance c.d.f., quantisation error and distortion}\label{S:distancecdt-etc}

%ZZZ TB rewritten
%Here we assume that $\SX$ is a connected compact subset of $\RR^d$ which coincides with the closure of its interior, as a cube, a ball, or a $d$-dimensional annulus. In Section~ZZZ we shall extend the results to the case where $\SX$ is the unit $d$-dimensional sphere $\SS_{d-1}(1)=\{\xb\in\RR^d: \, \|\xb\|=1\}$.

%$\mu$ is typically the uniform measure on $\SX$

%---------------------------------------------------------------------------
\subsection{Definitions and basic properties}

For any $n$-point set $\Xb_n$ we denote by $d(\cdot,\Xb_n)$ the distance function defined by $\xb\in\RR^d \mapsto d(\xb,\Xb_n)=\min_{\xb_i\in\Xb_n}\|\xb-\xb_i\|$; $\nu$ denotes a probability measure on $\RR^d$.

\begin{definition}\label{D:distance-cdf}
For $\Xb_n$ a fixed $n$-point set in $\RR^d$, the \emph{distance c.d.f.} $F(\cdot\,;\Xb_n,\nu)$ is the c.d.f.\ of the random variable $d(U,\Xb_n)$ when $U$ is distributed with $\nu$:
\bea %\label{F_Xn}
F(t;\Xb_n,\nu) = \nu\left\{U\in \bigcup_{i=1}^n \SB_d(\xb_i,t)\right\}  = \nu\{d(U,\Xb_n) \leq t\}\,.
\eea
When $\Xb_n=\Rb_n$ is a random $n$-point set generated with some probability measure $\Px$ on $\RR^{n\times d}$, we define the \emph{mean distance c.d.f.} $F_n(\cdot\,; \Px,\nu)$ of the random variable $d(U,\Rb_n)$ by $F_n(t; \Px,\nu) = \Ex_\Px\{F(t;\Rb_n,\nu) \}$.
\end{definition}

If $\nu$ is the delta measure $\delta_\ub$ concentrated at $\ub\in\RR^d$, the mean distance c.d.f.\ is simply $F_n(t; \Px,\delta_\ub)=\Px \left\{ d(\ub,\Rb_n)\leq t  \right\}$.

When $\SX$ is a compact subset of $\RR^d$, we denote by $\CR(\Xb_n)= \max_{\xb\in\SX} d(\xb,\Xb_n)$ the covering radius of $\Xb_n$. If $\nu$ is equivalent to the Lebesgue measure on $\SX$, $F(\cdot\,;\Xb_n,\nu)$ has a density on $[0,\CR(\Xb_n)]$, which we denote by $f(\cdot\,;\Xb_n,\nu)$, and the essential supremum of the random variable $d(U,\Xb_n)$ equals $\CR(\Xb_n)$.
When $\nu$ is the uniform measure on $\SX$, we simply denote the distance c.d.f.\ by $F(\cdot\,;\Xb_n)$: the random variable $d(U,\xb_n)$ is supported on $[0,\CR(\Xb_n)]$ and $F(\cdot\,;\Xb_n)$ contains all information about the filling of $\SX$ by $\Xb_n$. In particular, the $\mg$-quantiles $q_\mg(\Xb_n)$ of $F(\cdot\,;\Xb_n)$ provide useful space-filling characteristics. Assume that $\SX$ is connected and coincides with the closure of its interior. Then, the density $f(\cdot\,;\Xb_n)$ is strictly positive on $(0,\CR(\Xb_n))$ and $q_\mg(\Xb_n)$ is defined by
\bea %\label{eq:quant_eq}
F[q_\mg(\Xb_n);\Xb_n]=\mg \ \mbox{ for any }\mg\in(0,1)\,,
\eea
with $q_{0}(\Xb_n)=0$ and $q_{1}(\Xb_n)=\CR(\Xb_n)$.

\begin{definition}\label{D:quantisation-error}
The \emph{$L_s$-mean quantisation error} of a general probability measure $\mu$ on $\RR^d$ with finite $s$-th moment $\Ex\{\|X\|^s\}$ for a given $n$-point set $\Xb_n$ is the $L_s(\mu)$-norm of the distance function $d(\cdot,\Xb_n)$:
\bea %\label{E_s}
E_{\mu,s}(\Xb_n) = \|d(\cdot,\Xb_n)\|_{L_s} = \left[\int_\SX d^s(\xb,\Xb_n)\,\mu(\dd\xb)\right]^{1/s}\,, \quad s>0\,.\;\;\;
\eea
The quantity $E_{\mu,s}^s(\Xb_n)$ is called the $(\mu,s)$-distortion related to $\Xb_n$ {\rm \cite{Pages97}}. When $\Xb_n=\Rb_n$ is random, we define the expected $(\mu,s)$-distortion as $D_{\mu,s}(\Px)=\Ex_\Px\{E_{\mu,s}^s(\Rb_n)\}$.
\end{definition}

For any $s>0$, the $(\mu,s)$-distortion satisfies (see, e.g., \cite[p.~150]{Feller71-2}):
\be
E_{\mu,s}^s(\Xb_n) = \Ex_\mu\{d^s(U,\Xb_n)\} = s\, \int_0^\infty t^{s-1}\,[1-F(t;\Xb_n,\mu)] \,\dd t \,. \label{moment-ds-cdf}
\ee
The result is easily obtained when $\mu$ is such that $F(\cdot\,;\Xb_n,\mu)$ has a density $f(\cdot\,;\Xb_n,\mu)$. Indeed, by swapping the order of integration we get
\bea
\Ex_\mu\{d(U,\Xb_n)\} &=& \int_0^\infty \tau\,f(\tau;\Xb_n,\mu)\,\dd \tau = \int_0^\infty \left(\int_0^\tau \dd t\right) \,f(\tau;\Xb_n,\mu)\,\dd \tau \\
&& \hspace{-1cm} = \, \int_0^\infty \left(\int_t^\infty  f(\tau;\Xb_n,\mu)\,\dd \tau\right) \,\dd t = \int_0^\infty [1-F(t;\Xb_n,\mu)] \,\dd t \,.
\eea
Now, for any $s>0$, $G^{(s)}(\cdot\,;\Xb_n,\mu)$, defined by $G^{(s)}(t\,;\Xb_n,\mu)=F(t^{1/s};\Xb_n,\mu)$ for any $t\geq 0$, is the c.d.f.\ of $d^s(U,\Xb_n)$, and we obtain
\bea
\Ex_\mu\{d^s(U,\Xb_n)\} = \int_0^\infty [1-F(t^{1/s};\Xb_n,\mu)] \,\dd t = s\, \int_0^\infty t^{s-1}\,[1-F(t;\Xb_n,\mu)] \,\dd t \,. %\label{moment-ds-cdf}
\eea

%The expression \eqref{moment-ds-cdf} is valid more generally when $\mu$ is an arbitrary probability measure on $\RR^d$ (not necessarily absolutely continuous with respect to the Lebesgue measure); see \cite[p.~150]{Feller71-2}.  %and $\Xb_n$ is not necessarily included in $\SX$; see Section~ZZZ (annulus and sphere).

\begin{remark}\label{R:distortion-evaluation}
For any probability measure $\mu$ such that $\Ex\{\|X\|^{2s}\}<\infty$, the $(\mu,s)$-distortion of an arbirary $n$-point set $\Xb_n$ can be approximated by its Monte-Carlo estimator $E_{\mu_N,s}^s(\Xb_n)=(1/N)\sum_{i=1}^N d^s(U_i,\Xb_n)$ where the $U_i$ are i.i.d.\ with $\mu$, and $E_{\mu_N,s}^s(\Xb_n)$ satisfies a classical central limit theorem:
\bea
\sqrt{N} \left[E_{\mu_N,s}^s(\Xb_n)-E_{\mu,s}^s(\Xb_n)\right] \rad \SN(0,V(\Xb_n))\,, \quad N\to\infty\,,
\eea
where $V(\Xb_n)=E_{\mu,2s}^{2s}(\Xb_n)-E_{\mu,s}^{2s}(\Xb_n)$. When $\mu$ is supported on $\SX$ compact, $\CR(\Xb_n)<\infty$ and a direct application of Hoeffding's inequality gives:
\bea
\mbox{for all } \ma\in(0,1)\,, \ \mu\left\{|E_{\mu_N,s}^s(\Xb_n)-E_{\mu,s}^s(\Xb_n)|>\ma\, \CR(\Xb_n) \right\} < 2\,\e1^{-2\,N\ma^2} \,,
\eea
showing the exponentially fast concentration of $E_{\mu_N,s}^s(\Xb_n)$ to its mean $E_{\mu,s}^s(\Xb_n)$ as $N$ increases. When $\mu$ has unbounded support, depending on its tail properties, other concentration inequalities can also be derived.
\fin
\end{remark}


%---------------------------------------------------------------------------
\subsection{Quantisation error of random quantisers}

Let $\Rb_n=\{\xb_1,\ldots,\xb_n\}$ denote a random $n$-point set generated with a probability measure $\Px$ on $\RR^{n\times d}$. Our objective is to choose $\Px$ so that the expected $(\mu,s)$-distortion, $D_{\mu,s}(\Px)$ of Definition~\ref{D:quantisation-error}, is minimum.

From \eqref{moment-ds-cdf}, we can write
\be\label{Ls-mean-meandcdf}
D_{\mu,s}(\Px) =  s \int_{t\geq 0} t^{s-1}[1-F_n(t;\Px,\mu)]\,\dd t
\ee
where $F_n(\cdot\,; \Px,\mu) = \Ex_\Px\{F(t;\Rb_n,\mu) \}$ is the mean distance c.d.f.\ of Definition~\ref{D:distance-cdf}; that is, the c.d.f.\ of the random variable $d(U,\Rb_n)$ where both $U$ and $\Rb_n$ are random. The expected $(\mu,s)$-distortion $D_{\mu,s}(\Px)$ is the $(\mu,s)$-distortion for the c.d.f.\ $F_n(\cdot\,; \Px,\mu)$. In this paper we focus on distortion and quantisation error, but we could have also considered the $\mg$-quantiles of $F_n(\cdot\,; \Px,\mu)$; see Remark~\ref{R:quantiles}.
Note that Jensen inequality yields the following upper bound on the expected quantisation error: $\Ex_\Px\{ E_{\mu,s}(\Rb_n) \}\leq [D_{\mu,s}(\Px)]^{1/s}$ for $s \geq 1$.

By swapping the order of expectations, we get
\be
F_n(t; \Px,\mu)  &=& \Ex_\Px \left\{\Ex_\mu \left\{\Ind_{\{\ub\in\RR^d:\,d(\ub,\Rb_n)\leq t\}}(U) \right\}\right\} \nonumber \\
&=&  \Ex_\mu \left\{\Ex_\Px\left\{\Ind_{\{\ub\in\RR^d:\,d(\ub,\Rb_n)\leq t\}}(U) \right\}\right\} \nonumber \\
&=& \label{eq:prod_partial_cdf} \Ex_\mu  \left\{ F_n(t; \Px,\delta_U) \right\}\,,
\ee
where $F_n(t; \Px,\delta_\ub) = \Px  \left\{ d(\ub,\Rb_n)\leq t  \right\}$.

In the following, we consider random $n$-point set $\Rb_n$ with i.i.d.\ points having the distribution $\mathbb{P}$, so that $\Px(\dd\xb_1,\ldots,\dd\xb_n)=\prod_{i=1}^n \PP(\dd\xb_i)$; we shall denote this distribution $\Px=\PP^{[n]}$.
In Sections~\ref{S:PP-distribution-d2-Xspherical} and \ref{S:F-distribution-d2-Xspherical} we show that $F_n(t; \PP^{[n]},\mu)$ can be expressed in the form of a double integral when $\PP$ is spherically symmetric. This will rely on the following property which is true for any $\PP$.

For any fixed $\ub  \in \RR^d$, any $t \geq 0$ and any $\Rb_n$ with i.i.d.\ points having the arbitrary distribution $\PP$, we have
\be
F_n(t; \PP^{[n]},\delta_\ub) &=& 1-\prod_{j=1}^n \PP\left\{ \ub  \notin \SB_d({\xb}_j,t)  \right\} \nonumber \\
&& \hspace{-3cm} = \, 1-\prod_{j=1}^n\left(1- \PP  \left\{ \ub  \in \SB_d({\xb}_j,t)  \right\} \right) = 1-\bigg(1-\PP \left\{ \|X - \ub\| \leq t \right\} \bigg)^n\,,
\label{eq:prod_partial}
\ee
where $X \simd \PP$.

The expression \eqref{eq:prod_partial} indicates that the calculations of the mean distance c.d.f.\ $F_n(t; \PP^{[n]},\mu)$
%\eqref{eq:prod_partial_cdf}
and then of the expected $(\mu,s)$-distortion require the calculation of the probability
\bea %\label{eq:vol_inters}
\PP \left\{ \|X - \ub\| \leq t \right\} = \PP \left\{\ub    \in \SB_d(X,t)  \right\}\,,
\eea
which depends on $\PP$, $\ub $ and $r$. It is a consequence of the exchange of order of integration in \eqref{eq:prod_partial_cdf} that we have to consider another distance c.d.f., where now $\ub$ is fixed and plays the role of a one-point quantiser, whereas $X$ is random. Of course, this symmetry vanishes when considering random $n$-point sets, and the presence of $n$ i.i.d.\ points is accounted for explicitly in \eqref{eq:prod_partial}.
It is noticeable that this is the only place where the size $n$ appears.

%-----------------------------
\subsection{$\PP \left\{ \|X- \ub\| \leq t \right\}$ when $\PP$ is spherically symmetric}\label{S:PP-distribution-d2-Xspherical}

By definition, the random vector $X=(X_1, \ldots,X_d)$ is spherically symmetric if it can be represented as $X \eqd  R \cdot Z^{(d)} $, where $R=\|X\|$, $Z^{(d)}=(Z_1,\ldots,Z_d)$ is uniformly distributed on the unit sphere $\SS_{d-1}(1)$, and $R$ and $Z^{(d)}$ are independent.

When $\PP$ is spherically symmetric the distribution of $\|X  - \ub \|$ for $X\simd \PP$ only depends on $\|\ub\|$ and $d$. This distribution is specified in Proposition~\ref{th:dist} below. In the following, $\beta_{a,b}$ denotes a random variable with the Beta-density
\bea
\varphi_{a,b}(t)=t^{a-1}(1-t)^{b-1}/B(a,b) \ \mbox{ for } 0\leq t \leq 1\,,
\eea
where $B(a,b)$ is the Beta-function and $a,b>0$. The c.d.f.\ of $\beta_{a,b}$ is $\Pr\{\beta_{a,b} \leq t\}=I_t(a,b)$, where $I_t(a,b)$ is the regularised incomplete Beta-function, for which we use the convention
\be\label{convention}
I_t(\cdot,\cdot)=   \left \{\begin{array}{ll}
    0 &\;\;  {\rm for}\;\; t\leq 0 \\
  1 & \;\; {\rm for}\;\; t \geq 1 \, .\\
  \end{array}
\right.
\ee
The following lemma (see, e.g., \cite[Sect.~2.2]{FangKN90}) will be used several times.

\begin{lemma}\label{L:projections-beta}
Let $X=(X_1,\ldots,X_d) \eqd  R \cdot Z^{(d)}$ be a spherically symmetric random vector in $\RR^d$. Then, for any $m\in\{1,\ldots,d-1\}$, we have
\bea
X^{(m)}=(X_1,\ldots,X_m) \eqd  R \cdot \sqrt{\beta_{m/2,(d-m)/2}}\cdot Z^{(m)} \,,
\eea
where $R=\|X\|$, $\beta_{m/2,(d-m)/2}$ and $Z^{(m)}$ are independent.  Moreover, the joint density of $V^{(m)}=X^{(m)}/R$ is
\bea
\varphi_m(v_1,\ldots,v_m)= \frac{\Gamma(d/2)}{\pi^{m/2}\,\Gamma((d-m)/2)} \, \left(1-\sum_{i=1}^m v_i^2 \right)^{(d-m)/2-1} \ \mbox{ for } \sum_{i=1}^m v_i^2\leq 1\,.
\eea
\end{lemma}

The next proposition is the key element for the derivation of exact and asymptotic results on the behaviours of the mean distance c.d.f.\ and the $(\mu,s$)-distortion.

\begin{proposition} \label{th:dist}
Assume that $d\geq 2$ and $X\simd\PP$ with $\PP$ spherically symmetric. Then, for any fixed $\ub \in \RR^d$ we have
\be
 \|X  - \ub \|^2 \eqd (\|\ub\|-R)^2 +4\,\|\ub\|\,R \, \beta_{\delta,\delta}\, , \label{eq:quantile2}
\ee
where $\delta=(d-1)/2$ and the random variables $R=\|X \|$ and $\beta_{\delta,\delta}$ are independent.

Moreover, when $\Rb_n\sim \Px=\PP_a^{[n]}$ with $\PP_a$ uniform on $\SS_{d-1}(a)$, $a\geq 0$, we have
\be\label{d2URn}
d^2(\ub,\Rb_n)\eqd (\|\ub\|-a)^2+4\,a\,\|\ub\|\,\zeta(n,d)\,,
\ee
where $\zeta(n,d)=\min_{i=1,\ldots,n} \zeta_i$ and the $\zeta_i\eqd \beta_{\delta,\delta}$ are i.i.d.
\end{proposition}

\begin{proof}
As $X$ is spherically symmetric, the distribution of $\|X  - \ub \| $ only depends on $\ub $ through $r=\|\ub\|$, and without any loss of generality we can assume that $\ub=(r,0, \ldots, 0)$. We thus have
\bea %\label{eq:dist_sum}
 \|X \!-\! \ub \|^2
\eqd  (r\!-\!X_1)^2\!+\! \sum_{j=2}^d X_j^2 =  r^2\! -\! 2\,r X_1\!  +\!R^2= (R\!-\!r)^2 \!+\!2\,rR(1\!-\!Z_1)\,, \nonumber
\eea
where we have denoted $Z_1=X_1/R$. From Lemma~\ref{L:projections-beta}, $R$ and $Z_1$ are independent and the expression of $\varphi_1(\cdot)$ gives  $(1-Z_1)/2 \eqd  \beta_{\delta,\delta}$, which yields \eqref{eq:quantile2}.
%
Since the $n$ points of $\Rb_n$ are i.i.d.\ with $\PP_a$, the representation \eqref{eq:quantile2} yields \eqref{d2URn}.
%\carre
\end{proof}

From \eqref{eq:quantile2},  for all $ t\geq 0$ we have
\bea
\PP \left\{ \| X \!-\!\ub \|\! \leq \! { t } \right\}\!&=&  \! \Pr  \left\{(R-r)^2 +4\,Rr \, \beta_{\delta,\delta}\,  \leq  { t ^2} \right\}  \,,
\eea
with $r=\|\ub\| > 0$, and thus
\bea
 \PP  \left\{ \| X \!-\!\ub \|\! \leq \! { t } \right\}&=&\Pr  \left\{\, \beta_{\delta,\delta}\,  \leq  \frac{t^2 -(R-r)^2 }{4\,Rr}  \right\}\,. \;\;
\eea
By conditioning on the random variable $R= \|X \|$ and denoting $\Phi(\cdot )$ the c.d.f.\ of $R=\|X \|$, we obtain
\be\label{beta_approx_for_cube}
\PP  \left\{ \| X \!-\!\ub \|\! \leq \! { t } \right\}\!&=& \int_{\rho\geq 0}
\Pr  \left\{\, \beta_{\delta,\delta}\,  \leq  \frac{t^2 -(\rho-r)^2 }{4\,\rho r}  \right\} \dd\Phi(\rho )  \nonumber \\
&=& \int_{\rho\geq 0} I_\upsilon(\delta,\delta) \,\dd\Phi(\rho ) \,,
\ee
where
\be\label{v}
\upsilon=\upsilon(t,\rho,r)=[t^2 -(\rho-r)^2]/(4 \rho r)
\ee
and $I_\upsilon(\cdot,\cdot )$ is the regularised incomplete beta-function with the added convention~\eqref{convention}.
The integration in \eqref{beta_approx_for_cube} is over the support of the distribution of $R=\|X \|$.
If $r=\|\ub\| = 0$, then \eqref{beta_approx_for_cube} reduces to
\bea
    \PP  \left\{ \| X\!-\!\ub \|\! \leq \! { t } \right\} = \PP  \left\{\|X\| \leq t \right\} = \int_{0}^{{t}} \dd\Phi(\rho) \,.
\eea

When $\PP=\PP_a$, the uniform distribution on the sphere $\SS_{d-1}(a)$, the expression \eqref{beta_approx_for_cube} with the convention \eqref{convention} can be simplified as follows:
\be\label{PPaa}
\PP_a\left\{ \| X-\ub \| \leq t \right\} =
\left\{\begin{array}{ll}
0                      & \mbox{if } t\leq |r-a| \,,\\
I_\upsilon(\delta,\delta) & \mbox{if } |r-a| \leq t \leq a+r \,, \\
1                      & \mbox{if } a+r \leq t \,,
\end{array} \right.
\ee
where $r=\|\ub\|$ and $\upsilon=\upsilon(t,a,r)$ is given by \eqref{v}.

%-----------------------------
\subsection{$F_n(t; \PP^{[n]},\mu)$ and $D_{\mu,s}(\PP^{[n]})$ when $\PP$ is spherically symmetric}\label{S:F-distribution-d2-Xspherical}

From \eqref{eq:prod_partial_cdf} and \eqref{eq:prod_partial}, the mean distance c.d.f.\ $F_n(\cdot\,; \PP^{[n]},\mu)$ can be calculated explicitly as
\be \label{eq:prod_partial_cdf4}
F_n(t; \PP^{[n]},\mu) = 1- \Ex_\mu \left\{ \bigg(1-\PP \left\{ \| X-U \| \leq t \right\} \bigg)^n \right\}\,,
\ee
where $U\simd \mu$, $X \simd \PP$, and $X$ and $U$ are independent.
When $\PP$ is spherically symmetric, $\PP \left\{ \|X- \ub \| \leq t \right\} $ is given by \eqref{beta_approx_for_cube} which only depends on $\ub$ through $r=\|\ub\|$, and the expectation with respect to $U$ in \eqref{eq:prod_partial_cdf4} is reduced to an expectation with respect to the distribution of $\|U\|$. The calculation of $F_n(t; \PP^{[n]},\mu)$ thus amounts to a double integration, with respect to the distributions of $\|X\|$ and $\|U\|$.
From \eqref{Ls-mean-meandcdf}, the calculation of $D_{\mu,s}(\PP^{[n]})$ requires an additional integration with respect to $t$,
\be
D_{\mu,s}(\PP^{[n]}) &=&  s \int_{t\geq 0} t^{s-1} \Ex_\mu \left\{ \bigg(1-\PP \left\{ \| X-U \| \leq t \right\} \bigg)^n \right\} \,\dd t \,, \nonumber \\
&=& s \int_{t\geq 0} t^{s-1} \int_{r\geq 0} H(r,t;\Phi,n)\, \dd\Psi(r)\,\dd t \,, \label{Ls-mean-final}
\ee
where
\bea %\label{H}
H(r,t;\Phi,n) =\bigg(1-\int_{\rho\geq 0} I_\upsilon(\delta,\delta) \,\dd\Phi(\rho ) \bigg)^n \,,
\eea
with $\upsilon=\upsilon(t,\rho,r)$ given by \eqref{v}, $\Psi(\cdot)$ the c.d.f.\ of $\|U\|$ for $U\simd\mu$ and $\Phi(\cdot)$ the c.d.f.\ of $\|X\|$ for $X\simd\PP$.

The integrals required to compute $D_{\mu,s}(\PP^{[n]})$ can be evaluated with arbitrary precision. When $\PP$ is appropriately parameterised, a high-performance random quantiser can be obtained by numerically minimising the distortion with respect to its defining parameters. For the remainder of this paper, we focus on quantisers defined by a distribution $\PP$ that depends on a single scalar parameter: the radius of a sphere (Section~\ref{S:sphere}), the radius of a ball (Section~\ref{S:ball}), or the variance of a multivariate normal distribution (Section~\ref{S:normal}). The optimisation is performed using a derivative-free line-search method, specifically the golden-section algorithm.

%where $\PP \left\{ \| X-U \| \leq t \right\}$ is calculated via \eqref{beta_approx_for_cube}.
%
%When $\PP=\PP_a$ uniform on $\SS_{d-1}(a)$, and \eqref{Ls-mean-final} gives
%\bea
%D_{\mu,s}(\PP_a^{[n]}) =
%\eea
%given by \eqref{PPaa}.


\begin{remark}
The expression \eqref{Ls-mean-final} gives the expected $(\mu,s)$-distortion for random quantisers. Therefore, according to the paradigm of the so-called ``probabilistic method" (see, e.g., \cite{AlonS2000}), for any $\PP$ there exists at least one non-random $n$-point set $\Xb_n$ with $(\mu,s)$-distortion $E_{\mu,s}^s(\Xb_n)\leq D_{\mu,s}(\PP^{[n]})$. See Section~\ref{S:sphere-2^d} for an illustration involving full factorial designs.
\fin
\end{remark}

%{\Luc ZZZ TB rewritten

%In the next two sections we apply the results above to the construction of quantisers in the case where $\SX$ is an annulus (with the ball and the sphere as particular instances) or a cube. In the annulus, although $\mu$ can be an arbitrary spherically symmetric probability measure, we restrict our attention to the case where $\mu$ is uniform in $\SX$. ZZZ try $\mu$ = normal? In the cube, we shall also consider $\mu$ uniform on $\SX$, which is not spherically symmetric. However, suitable approximations will be derived that allow us to extend the results above.}

%___________________________________________________________________________
\section{Exact, non-asymptotic, results}\label{S:ball-sphere-annulus}

Throughout this section, three emblematic cases of  spherically symmetric measures are investigated. Nevertheless, the results remain valid for any measure $\nu$, regardless of its symmetr, so long as the radial component $r=\|U\|$ follows the same distribution under $\nu$ as under $\mu$. For instance, in Section~\ref{S:sphere} ($\mu$ uniform on $\SS_{d-1}(1)$), $\nu$ can be a point mass $\delta_\ub$ for any unit vector $\ub$; in Sections~\ref{S:ball} ($\mu$ uniform in $\SB_d(1)$) and~\ref{S:normal} ($\mu$ normal), $\nu$ may be supported on a line segment or a line, respectively, provided the distribution of $r$ is preserved.

%-----------------------------
\subsection{$\mu$ is uniform on the unit sphere $\SS_{d-1}(1)$} \label{S:sphere}

Here $\mu$ is the spherical measure (uniform on $\SS_{d-1}(1)$), and the distribution of $\|U\|$ for $U\simd\mu$ is the delta measure at $r=1$. When $n=1$ and $\Xb_n=\{\0b_d\}$, the distance c.d.f.\ $F(t;\Xb_n,\mu)$ equals 0 for $t<1$ and 1 for $t\geq 1$. For any other $\Xb_n$ it has a density $f(\cdot\,;\Xb_n,\mu)$ on $[0,\CR(\Xb_n)]$.
%and we can apply the developments of Section~\ref{S:distancecdt-etc}.

We first consider random quantisers with distribution $\PP=\PP_a$ uniform on $\SS_{d-1}(a)$  and minimise $D_{\mu,s}(\PP^{[n]}_a)$ with respect to $a\in[0,1]$. The expression \eqref{PPaa} can be simplified as follows:
\bea %\label{PPa}
\PP_a\left\{ \| X-\ub \| \leq t \right\} =
\left\{\begin{array}{ll}
0                      & \mbox{if } t\leq 1-a \,,\\
I_\upsilon(\delta,\delta) & \mbox{if } 1-a \leq t \leq 1+a \,, \\
1                      & \mbox{if } 1+a \leq t \,,
\end{array} \right.
\eea
where $\upsilon=[t^2 -(1-a)^2]/(4 a)$. Formula~\eqref{Ls-mean-final} then gives
\be\label{Ds-sphere}
D_{\mu,s}(\PP^{[n]}_a) =  (1-a)^s + s \int_{1-a}^{1+a} t^{s-1} \left[1-I_\upsilon(\delta,\delta)\right]^n\,\dd t \,.
\ee
This exact expression for the expected $(\mu,s)$-distortion depends on $d,n,s$ and $a$; the value $a^*=a^*(d,n,s)$ that minimises $D_{\mu,s}(\PP^{[n]}_a)$ with respect to $a$ can be obtained numerically with arbitrary precision for any given $d,n$ and $s$. In the special case $d=3$, $\delta=1$ and $I_v(1,1)=v$ for $v\in[0,1]$, so that $D_{\mu,s}(\PP^{[n]}_a)$ can be calculated explicitly for $s$ even. It is given by a polynomial in $a$ of degree $s$, with for example $D_{\mu,2}(\PP^{[n]}_a)=(1-a)^2+4\,a/(n+1)$ and $D_{\mu,4}(\PP^{[n]}_a)=(1-a)^4+8\,a[n(1-a)^2+2+2\,a^2]/[(n+1)(n+2)]$, which gives $a^*(3,n,2)=(n-1)/(n+1)$ and $a^*(3,n,4)$ as the root of a cubic equation, which satisfies $a^*(3,n,4)=1-4/n+16/n^2+\SO(1/n^3)$ as $n\to\infty$.

The left column of Figure~\ref{F:sphere} presents $a^*(d,n,s)$ as a function of $d$ for different $n$ and $s=1$ (top) and $s=10$ (bottom). For fixed $s$ and $d$, $a^*$ increases with $n$, and asymptotically, when $n$ tends to infinity, $a^*$ tends to 1 (see, e.g., \cite[Chap.~9]{GrafL2000}). Intuitively the choice $a=1$ seems natural for all $n$ and $s$. However, Figure~\ref{F:sphere} indicates that $a=1$ can be a poor choice, especially in high dimension. For any fixed $d$ and $n$, $a^*(d,n,s) \to 1$ when $s\to 0$ and is a decreasing function of $s$, with $a^*(d,n,\infty)=0$ if $n$ is small enough. Indeed, for $n \leq d$, $\CR(\Xb_n)\geq 1$ for any $\Xb_n$ and the inequality is strict if $\0b_n\not\in\Xb_n$. This follows from the fact that any $d$ points in $\RR^d$ belong to a $d$-dimensional hyperplane $\SH_d$; one of the two parts of $\SS_{d-1}(1)$ cut by $\SH_d$ necessarily contains a hemisphere whose pole is at distance at least one from each point. This implies that $a^*(d,n,s)=0$ for large enough $s$ when $n\leq d$. Also, for fixed $n$ and $s$, $a^*(d,n,s)$ decreases with $d$.

The central column of Figure~\ref{F:sphere} shows that, as expected, $D_{\mu,s}^{1/s}(\PP^{[n]}_{a^*})$ decreases with $n$ and increases with $d$ and $s$. The right column presents the efficiency of the naive quantiser with $n$ points i.i.d.\ on $\SS_{d-1}(1)$ relative to optimised random quantisers with $n$ points on $\SS_{d-1}(a^*)$. It shows that the benefit of using an appropriate $a$ is significant when $d$ is large or $n$ is small.

\begin{figure}[ht!]
\begin{center}
\includegraphics[width=.32\linewidth]{Rstar_sphere_s1_D_n}
\includegraphics[width=.32\linewidth]{Qstar1s_sphere_s1_D_n}
\includegraphics[width=.32\linewidth]{Qstar1s-over-Q11s_sphere_s1_D_n} \\
\includegraphics[width=.32\linewidth]{Rstar_sphere_s10_D_n}
\includegraphics[width=.32\linewidth]{Qstar1s_sphere_s10_D_n}
\includegraphics[width=.32\linewidth]{Qstar1s-over-Q11s_sphere_s10_D_n}
\end{center}
\caption{\small Sphere. Value $a^*$ minimising $D_{\mu,s}(\PP^{[n]}_a)$ (left column), $D_{\mu,s}^{1/s}(\PP^{[n]}_{a^*})$ (central column) and ratio $D_{\mu,s}^{1/s}(\PP^{[n]}_{a^*})/D_{\mu,s}^{1/s}(\PP^{[n]}_1)$ (right column) as functions of $d$ for $s=1$ (top row) and $s=10$ (bottom row) and different values of $n$: $n=10$ ({\color{Cerulean} $\bullet$}), $n=10^2$ ({\color{Dandelion} $\blacklozenge$}), $n=10^3$ ({\color{green} $\blacktriangledown$}), $n=10^4$ ({\color{Orchid} $\bigstar$}), and $n=10^5$ ({\color{gray} {\tiny $\blacksquare$}}).}
\label{F:sphere}
\end{figure}
% tst_quantization_random_sphere.m

\vsp
The next figure (Figure~\ref{F:sphere2}, where $s=10$, $d=10$) shows that adding a delta measure at $\0b_d$ to $\PP_1$ may also yield a decrease of the expected distortion for small~$n$. Denote $\QQ_{\ma,a}=\ma\,\PP_a+(1-\ma)\,\delta_{\0b_d}$. Similarly to \eqref{Ds-sphere}, we obtain
\bea
D_{\mu,s}(\QQ^{[n]}_{\ma,a})  &=&  (1-a)^s + s \int_{1-a}^1 t^{s-1} \left[1-\ma\, I_\upsilon(\delta,\delta)\right]^n\,\dd t \\
&& + s\,\ma^n \int_1^{1+a} t^{s-1} \left[1- I_\upsilon(\delta,\delta)\right]^n\,\dd t \,,
\eea
with $\upsilon=[t^2 -(1-a)^2]/(4 a)$.
The left panel presents the optimal $\ma^*$ minimising $D_{\mu,s}(\QQ^{[n]}_{\ma,1})$ as a function of $n$: unsurprisingly, $\ma^*$ quickly approaches 1 as $n$ increases. One may observe that, for all $n$, $(1-\ma^*)>1/n$ (it corresponds to the probability that a realisation of a random set $\Rb_n$ contains a point at $\0b_d$).
Next, for a given $n$ ($n=20$), we use the corresponding optimal $\ma^*$ minimising $D_{\mu,s}(\QQ^{[n]}_{\ma,1})$ ($\ma^*\simeq 0.846$) and plot $D_{\mu,s}^{1/s}(\QQ^{[n]}_{\ma^*,a})$ and $D_{\mu,s}^{1/s}(\QQ^{[n]}_{1,a})=D_{\mu,s}^{1/s}(\PP^{[n]}_a)$ as functions of $a$ (central panel): the effect of introducing a delta measure at $\0b_d$ is significant for $a$ close to one, but choosing a suitable radius $a$ has bigger impact. Finally, we consider the efficiencies of random quantisers with $\QQ_{\ma^*,1}$ with respect to quantisers distributed with $\PP_1$ and
$\PP_{a^*}$ (with $\ma^*$ and $a^*$ depending on $n$) as functions of $n$ (right panel): when $a$ is set to 1, the choice of a suitable $\ma$ has some positive effect for small $n$, but the impact of choosing a suitable $a$ is stronger for all $n$.

\begin{figure}[ht!]
\begin{center}
\includegraphics[width=.32\linewidth]{Qstar1s_sphere_s10-d10_alphastar1_n}
\includegraphics[width=.32\linewidth]{Qstar1s_sphere_s10-d10_alpha1-alphastar1_a}
\includegraphics[width=.32\linewidth]{Qstar1s_sphere_s10-d10_ratios-n} \\
\end{center}
\caption{\small Left: value $\ma^*$ minimising $D_{\mu,s}(\QQ^{[n]}_{\ma,1})$ as function of $n$. Center: $D_{\mu,s}^{1/s}(\QQ^{[n]}_{\ma^*,a})$ ({\color{red} {\bf - - -}}) and $D_{\mu,s}^{1/s}(\QQ^{[n]}_{1,a})$ ({\color{Cerulean} {\bf ---}}) as functions of $a$ for $n=20$ and $\ma^* = 0.846$. Right: $D_{\mu,s}^{1/s}(\QQ^{[n]}_{\ma^*,1})/D_{\mu,s}^{1/s}(\PP^{[n]}_1)$ ({\color{Cerulean} $\bullet$}) and $D_{\mu,s}^{1/s}(\QQ^{[n]}_{\ma^*,1})/D_{\mu,s}^{1/s}(\PP^{[n]}_{a^*})$ ({\color{Dandelion} $\blacklozenge$}) as functions of $n$; $s=10$, $d=10$.}
\label{F:sphere2}
\end{figure}
% tst_quantization_random_sphere.m

%-----------------------------
\subsection{Comparison between random quantisers uniform on a sphere and optimised full factorial $2^d$ designs} \label{S:sphere-2^d}
Let $n=2^d$ and $\Xb_n(b)$ be the full factorial $2^d$ design with points $\xb_i$ having coordinates $\pm b$, $b\in(0,1]$: the $\xb_i$ are vertices of the cube $[-b,b]^d$; they also belong to the sphere $\SS_{d-1}(a)$ with $a=\sqrt{d}\,b$. For $\SX=\SS_{d-1}(1)$, the Voronoi regions $\SV(\xb_i)=\{\xb\in\SX: \|\xb-\xb_i\| = \min_{\xb_j\in\Xb_n(b)} \|\xb-\xb_j\|\}$ are all identical up to a permutation of coordinates and do not depend on the value of $b$.

Consider first the case $s=2$. For $\mu$ uniform on $\SX=\SS_{d-1}(1)$, $E_{2,\mu}^2[\Xb_n(b)]=\Ex\{\|V-\bb\|^2\}$ where $\bb=(b,\ldots,b)\in\RR^d$ and $V=(V_1,\ldots,V_d)$ is uniform in $\SV(\bb)=\SS_{d-1}(1)\cap \RR_+^d$. Therefore
\bea
E_{2,\mu}^2[\Xb_n(b)]= \Ex\{ \|V\|^2 - 2\,b\sum_{i=1}^d  V_i + d b^2\} = 1 - 2\,bd\, \Ex\{V_1\} + db^2 \,.
\eea
Minimisation of $E_{2,\mu}^2[\Xb_n(b)]$ with respect to $b$ yields $b_2^*=\Ex\{V_1\}$ and the optimal $(\mu,2)$-distortion of a $2^d$ factorial design is $E_{2,\mu}^2[\Xb_n(b_2^*)]=1-d (b_2^*)^2$.
For $U^{(d)}=(U_1,\ldots,U_d)\simd \mu$, the first component $U_1$ has the density $\varphi_1(\cdot)$ of Lemma~\ref{L:projections-beta} on $[-1,1]$, therefore, $V_1$ has the density $2\,\varphi_1(\cdot)$ on $[0,1]$, so that
\bea
b_2^* =\Ex\{V_1\} &=& \frac{\Gamma(d/2)}{\sqrt{\pi}\,\Gamma((d+1)/2)}\,, \\
E_{2,\mu}^2[\Xb_n(b_2^*)] &=& 1 - \frac{d}{\pi} \left(\frac{\Gamma(d/2)}{\Gamma((d+1)/2)}\right)^2 \,.
\eea

More generally, the analytic expression for the $(\mu,s)$-distortion $E_{\mu,s}^s[\Xb_n(b)]$ can be derived for any even $s$, but the calculations become more complicated as $s$ increases. For $s=4$ we obtain
\bea
E_{4,\mu}^4[\Xb_n(b)] &=& \Ex\{\|V-\bb\|^4\} \\
&=& 1 +4\,b^2\, \Ex\left\{\left(\sum_{i=1}^d V_i \right)^2\right\} + d^2 b^4 + 2\,d b^2 - 4\, bd\, \Ex\{V_1\} (1+d b^2) \,.
\eea
Using $\Ex\{(\sum_{i=1}^d V_i)^2\}=1+d(d-1) \Ex\{V_1V_2\}$, where $\Ex\{V_1V_2\}=2/(\pi d)$ from $\varphi_2(\cdot,\cdot)$ of Lemma~\ref{L:projections-beta}, together with the expression above for $\Ex\{V_1\}$, we can express $E_{4,\mu}^4[\Xb_n(b)]$ as a fourth-degree polynomial in $b$ with coefficients depending on $d$. The optimal value $b_4^*$ can be obtained explicitly.
The values of $b_2^*$ and $b_4^*$ are very close for large $d$ as $a_j^*=\sqrt{d}\,b_j^*=\sqrt{2/\pi}+\SO(1/d)$ for $j=2,4$.

For $s=\infty$, $E_{\infty,\mu}[\Xb_n(b_\infty^*)]=\CR[\Xb_n(b)]$, the covering radius of $\Xb_n(b)$, with $(\CR[\Xb_n(b)])^2=1+d\,b^2-2\,b$ being minimum for $b_\infty^*=1/d$, which gives
$\CR[\Xb_n(b_\infty^*)]=\sqrt{1-1/d}$.

Figure~\ref{F:sphere3} compares the optimal full factorial designs $\Xb_n(b_s^*)$ for $s=2$ and 4 with optimised random quantisers $\PP_{a^*}^{[n]}$ on $\SS_{d-1}(1)$ with the same sample size $n=2^d$. The left panel shows the values of the radii of the various spheres on which the points lie, as functions of $d$. The right panel presents $D_{\mu,s}^{1/s}(\PP^{[n]}_{a^*})$ for the random quantisers and $E_{\mu,s}[\Xb_n(b_s^*)]$ for the full factorial designs, for $s=2$ and 4. For both types of point sets, the behaviours for $s=2$ and $s=4$ are similar. As $d$ tends to infinity, for the full factorial designs $a_s^*=\sqrt{d}\,b_s^*$ tends to the limit $\sqrt{2/\pi}\simeq 0.797885$ from above. For random quantisers, we observe that $a^*(s)=a^*(d,n,s)$ defined in Section~\ref{S:sphere} tends from below to a larger limiting value---we shall see in Section~\ref{S:extreme-value-sphere} that this limiting value equals $\sqrt{3}/2 \simeq 0.8660$.
The expected $(\mu,s)$-distortion of random quantisers shows little sensitivity to the choice of $a$ in the neighborhood of $a^*(s)$.
In particular, the plots of $D_{\mu,s}^{1/s}(\PP^{[n]}_{a_s^*})$ (not shown) and $D_{\mu,s}^{1/s}(\PP^{[n]}_{a^*(s)})$ almost coincide for $d>3$.

For both values of $s$, $a_s^* \approx a^*(s)$ for $d=7$. This dimension is the critical one at which random quantisers start exhibiting a smaller quantisation error than full factorial designs. As a consequence of the probabilistic method, it implies that for all $d\geq 7$ there exist non-random quantisers with $n=2^d$ that have smaller quantisation errors for $s=2$ and 4 than optimised full factorial designs.


\begin{figure}[ht!]
\begin{center}
\includegraphics[width=.49\linewidth]{astar_sphere_s2-s4_factorial}
\includegraphics[width=.49\linewidth]{Esmu_sphere_s2-s4_factorial} \\
\end{center}
\caption{\small Left: Values $a^*(s)$ minimising $D_{\mu,s}(\PP^{[n]}_a)$ for $n=2^d$ with $s=2$ ({\color{Dandelion} $\blacklozenge$}) and $s=4$ ({\color{red} $\bigstar$}), and values $a_s^*=\sqrt{d}\,b_s^*$ minimising $E_{\mu,s}[\Xb_n(b)]$ for $s=2$ ({\color{Cerulean} $\bullet$}) and $s=4$ ({\color{blue} {\tiny $\blacksquare$}}), as functions of $d$.
Right: $D_{\mu,s}^{1/s}(\PP^{[n]}_{a^*})$ for $n=2^d$ with $a^*=a^*(d,n,s)$, $s=2$ ({\color{Dandelion} $\blacklozenge$}) and $s=4$ ({\color{red} $\bigstar$}), and values $E_{\mu,s}[\Xb_n(b_s^*)]$ for $s=2$ ({\color{Cerulean} $\bullet$}) and $s=4$ ({\color{blue} {\tiny $\blacksquare$}}), as functions of $d$.}
\label{F:sphere3}
\end{figure}
% tst_quantization_random_sphere.m

%-----------------------------
\subsection{$\mu$ is uniform in the unit ball $\SB_d(1)$} % or an annulus}
\label{S:ball}

%-----------------------------
%\subsection{$\SX$ is an annulus: $\SX=\SA_{\me,1}=\SB_d(1)\setminus\SB_d^o(\me)$} \label{S:annulus}

%For any $\me\in[0,1)$, we consider the annulus $\SA_d(\me,1)=\SB_d(1)\setminus\SB_d^o(\me)$: in the particular case $\me=0$, $\SA_d(0,1)$ is the ball $\SB_d(1)$; when $\me=1$ we have $\SA_d(1,1)=\SS_{d-1}(1)$, the situation considered in Section~\ref{S:sphere}.

%For $\mu$ uniform in $\SA_d(\me,1)$, the density of $r=\|U\|$ for $U\simd\mu$ is $\psi_\me(\tau)=[d/(1-\me^d)]\,\tau^{d-1}$, $\me\leq \tau \leq 1$. Note that for any given $\me$ the densities $\psi_\me(\cdot)$ and $\psi_0(\cdot)$ become very close for large $d$, since $\psi_0(\tau)/\psi_\me(\tau)=1-\me^d$ for $\tau\in[\me,1]$.

For $\mu$ uniform in $\SB_d(1)$, the density of $r=\|U\|$ for $U\simd\mu$ is $\psi(\tau)=d\,\tau^{d-1}$.

%We consider random quantisers whose $n$ points are i.i.d.\ with $\PP=\PP_{a,b}$ uniform on an annulus $\SA_d(a,b)=\SB_d(b)\setminus\SB_d^o(a)$, where the tuning parameters $a$ and $b$ satisfy $0\leq a \leq b \leq 1$.

%When $X\simd\PP_{a,b}$, the c.d.f.\ $\Phi(\cdot)$ of $R=\|X\|$ has the density $\phi_{a,b}(\rho)=[d/(b^d-a^d)]\,\rho^{d-1}$, $\rho\in[a,b]$.

%{\Luc
%The sub-Gaussian norm $\|R\|_{\psi_2}$ of $R$, defined by $\|R\|_{\psi_2}=\inf\{t>0: \, \Ex\{\e1^{R^2/t^2}\} \leq 2\}$ is smaller than $b/\sqrt{\log 2}$
%Since the random variable $R$ is bounded, it is sub-Gaussian. Its variance proxy equals $(b-a)^2/4$ and its sub-Gaussian norm $\|R\|_{\psi_2}=\inf\{t>0: \, \Ex\{\e1^{R^2/t^2}\} \leq 2\}$ is smaller than $b/\sqrt{\log 2}$; see \cite[Sect.3]{ZhangC2021}. The condition \eqref{norm-concentration} of Proposition~\ref{Pr:norm-concentration} is therefore satisfied.
%}

%In the limiting case $a=b$, $\PP_{a,a}=\PP_a$ is uniform on $\SS_{d-1}(a)$, and $\PP_a\{ \| X-\ub \| \leq t \}$ is given by \eqref{PPaa}. We consider this case first.

Consider first the case where $\PP=\PP_a$ uniform on $\SS_{d-1}(a)$: \eqref{Ls-mean-final} gives
\be\label{DS-Pa}
D_{\mu,s}(\PP_a^{[n]}) =  s \int_{t\geq 0} t^{s-1} \int_{r\geq 0}  H(r,t;\Phi_a,n) \,\psi(r)\,\dd r \,\dd t \,,
\ee
where $H(r,t;\Phi_a,n) =\bigg(1-\PP_a \left\{ \| X-\ub \| \leq t \right\} \bigg)^n$ with $\Phi_a(\cdot)$ the c.d.f.\ of the delta measure $\delta_a$ and $r=\|\ub\|$, so that $\PP_a\{ \| X-\ub \| \leq t \}$ is given by \eqref{PPaa}.

We also consider random quantisers whose $n$ points are i.i.d.\ with $\PP=\PP_{0,b}$ uniform in $\SB_d(b)$.
The c.d.f.\ $\Phi(\cdot)$ of $R=\|X\|$ has then the density $\phi_{0,b}(\rho)=d\,\rho^{d-1}/b^d$ for $\rho\in[0,b]$.
More general spherically symmetric distributions $\PP$ could also be used, but we found that, unless $n$ is extremely small, the class of distributions $\PP_{0,b}$ is rich enough.
The expected $(\mu,s)$-distortion $D_{\mu,s}(\PP_{0,b}^{[n]})$ is given by \eqref{Ls-mean-final}
where
\bea
H(r,t;\Phi,n) =\bigg(1- \int_a^b I_\upsilon(\delta,\delta) \,\phi_{0,b}(\rho)\,\dd\rho \bigg)^n \,,
\eea
with $\upsilon=\upsilon(t,\rho,r)$ given by \eqref{v}.

Figure~\ref{F:ball} presents the same information as Figure~\ref{F:sphere} for the ball configuration: we consider optimal design defined by $\PP_{0,b^*}$, where $b^*=b^*(d,n,s)$ minimises $D_{\mu,s}(\PP_{0,b}^{[n]})$ and the right column compares performances of $\PP_{0,b^*}$ and $\PP_{0,1}$.

\begin{figure}[ht!]
\begin{center}
\includegraphics[width=.32\linewidth]{Rstar_ball_s1_D_n}
\includegraphics[width=.32\linewidth]{Qstar1s_ball_s1_D_n}
\includegraphics[width=.32\linewidth]{Qstar1s-over-Q11s_ball_s1_D_n} \\
\includegraphics[width=.32\linewidth]{Rstar_ball_s10_D_n}
\includegraphics[width=.32\linewidth]{Qstar1s_ball_s10_D_n}
\includegraphics[width=.32\linewidth]{Qstar1s-over-Q11s_ball_s10_D_n}
\end{center}
\caption{\small Ball. Value $b^*$ minimising $D_{\mu,s}(\PP^{[n]}_{0,b})$ (left column), $D_{\mu,s}^{1/s}(\PP^{[n]}_{0,b^*})$ (central column) and ratio $D_{\mu,s}^{1/s}(\PP^{[n]}_{0,b^*})/D_{\mu,s}^{1/s}(\PP^{[n]}_{0,1})$ (right column) as functions of $d$ for $s=1$ (top row) and $s=10$ (bottom row) and different values of $n$: $n=10$ ({\color{Cerulean} $\bullet$}), $n=10^2$ ({\color{Dandelion} $\blacklozenge$}), $n=10^3$ ({\color{green} $\blacktriangledown$}), $n=10^4$ ({\color{Orchid} $\bigstar$}), and $n=10^5$ ({\color{gray} {\tiny $\blacksquare$}}).}
\label{F:ball}
\end{figure}
% tst_quantization_random_ball.m

The quantities considered behave qualitatively as in Figure~\ref{F:sphere}.  Note, however, that $b^*$ is larger than $a^*$ plotted in Figure~\ref{F:sphere} (left columns---$b^*$ is much larger than $a^*$ on the bottom row with $s=10$), that $D_{\mu,s}(\PP_{0,b^*}^{[n]})$ is smaller than $D_{\mu,s}(\PP_{a^*}^{[n]})$ (central columns), and that the efficiency of a ``classical" random design compared to an optimised design is slightly larger for the ball than for the sphere (right columns). %The central and left columns thus suggests that quantisation is easier for $\mu$ uniform in $\SB_d(1)$ than for $\mu$ uniform in $\SS_{s-1}(1)$.
For a fixed dimension $d$, the empirical distribution of an optimal quantiser which minimises $E_{\mu,s}(\Xb_n)$ converges weakly to $\PP_{0,1}^{[n]}$ as $n \to \infty$ (see \cite[Table~7.1]{GrafL2000}). Consistent with this theoretical result, we empirically observe that the optimal parameters $b^*$ and $a^*$ both converge to 1 as $n \to \infty$; see also Section~\ref{S:extreme-value-ball} for further details on their asymptotic behaviour.

%Figure~\ref{F:a_star-bstar_ball}
Figure~\ref{F:bstar_ALLastar_ball_s2_D} of Section~\ref{S:extreme-value-ball} displays the optimal values $b^*(d,n,s)$ and $a^*(d,n,s)$, respectively minimising $D_{\mu,s}(\PP_{0,b}^{[n]})$ and $D_{\mu,s}(\PP_a^{[n]})$, plotted as functions of $d$ for $s=2$ and two distinct values of $n$. The figure shows that $b^*>a^*$, and the gap $b^*-a^*$ narrows as $d$ increases, reflecting the fact that both $a^*$ and $b^*$ converge to zero as $d\to\infty$. The behaviour remains qualitatively unchanged for other values of $s$, especially when the dimension $d$ is large.

%\begin{figure}[ht!]
%\begin{center}
%\includegraphics[width=.49\linewidth]{bstar_astar_ball_s10_n1000_D}
%\includegraphics[width=.49\linewidth]{bstar_astar_ball_s10_n100000_D}
%\end{center}
%\caption{\small Optimal values $a^*(d,n,10)$ ({\color{Cerulean} $\bullet$}) and $b^*(d,n,10)$ ({\color{Dandelion} $\blacklozenge$}) as functions of $d$ for $n=1\,000$ (left) and $n=100\,000$ (right).}
%\label{F:a_star-bstar_ball}
%\end{figure}
% tst_quantization_random_ball.m

Given the observations in Figures~\ref{F:ball} and %\ref{F:a_star-bstar_ball}
\ref{F:bstar_ALLastar_ball_s2_D}, two key patterns emerge:
(\textit{i}) for large dimensions $d$, the values of $a^*(d,n,s)$ and $b^*(d,n,s)$ become very close; (\textit{ii}) the quantiser $\PP_{0,b^*}^{[n]}$ significantly outperforms $\PP_{0,1}^{[n]}$ when $d$ is sufficiently large. Moreover, we also observe that $\PP_{a^*}^{[n]}$ outperforms $\PP_{0,b^*}^{[n]}$ for reasonable design sizes $n$; see Figure~\ref{F:ALLeff_ball_s2_D} in Section~\ref{S:extreme-value-ball}.
Since $b^*$ converges to 1 as $n \to \infty$, with $\PP_{0,1}^{[n]}$ being asymptotically optimal, a natural question arises: for a fixed dimension $d$, at what sample size $n=n^*(d,s)$ does $\PP_{a^*}^{[n]}$ stop being preferable to $\PP_{0,b^*}^{[n]}$? Empirical analysis reveals that this transition occurs only for extremely large values of $n$.
Define
\be\label{nstar-ball}
n^*(d,s) = \max\{n\in\NN: D_{\mu,s}(\PP_{a^*}^{[n]}) < D_{\mu,s}(\PP_{0,b^*}^{[n]}) \} \,.
\ee
The left panel of Figure~\ref{F:nstar-ball_s2-4-10} presents $n^*(d,s)$ (in log scale, obtained by dichotomous search) as a function of $d=3,\ldots,20$ for three values of $s$, indicating a superexponential increase of $n^*(d,s)$ with $d$.

\begin{figure}[ht!]
\begin{center}
\includegraphics[width=.49\linewidth]{nstar-ball_d3-20_s2-4-10}
\includegraphics[width=.49\linewidth]{nstar-normal_d3-20_s2-4-10}
\end{center}
\caption{\small $n^*(d,s)$ as a function of $d$ for $s=2$ ({\color{Dandelion} $\blacklozenge$}), $s=4$ ({\color{red} $\bigstar$}) and $s=10$ ({\color{Cerulean} $\bullet$}). Left: $\mu$ is uniform on $\SS_{d-1}(1)$ and $n^*(d,s)$ is given by \eqref{nstar-ball}. Right: $\mu$ is spherically normal and $n^*(d,s)$ is given by \eqref{nstar}.
}
\label{F:nstar-ball_s2-4-10}
\end{figure}
% tst_quantization_random_ball.m

%\vsp
%Consider now the case where $\PP=\PP_{a,b}$ is uniform on the annulus $\SA_d(a,b)$. The density of $R=\|X\|$ with $X\simd\PP_{a,b}$ is $\phi_{a,b}(\rho)=[d/(b^d-a^d)]\,\rho^{d-1}$, $\rho\in[a,b]$. The expected $(\mu,s)$-distortion $D_{\mu,s}(\PP_a^{[n]})$ is given by \eqref{Ls-mean-final} where
%\bea
%H(r,t;\Phi,n) =\bigg(1- \int_a^b I_\upsilon(\delta,\delta) \,\phi_{a,b}(\rho)\,\dd\rho \bigg)^n \,,
%\eea
%with $\upsilon=\upsilon(t,\rho,r)$ given by \eqref{v}.

%-----------------------------
\subsection{$\mu$ is a spherically symmetric normal distribution}\label{S:normal}

Any spherically symmetric normal distribution in $\RR^d$ can be renormalised as $\SN(\0b_d,\Ib_d/d)$, with $\Ib_d$ the $d$-dimensional identity matrix. In this section, we assume that $\mu$ is the corresponding probability measure. When $U\simd\mu$, $d\,\|U\|^2$ has the chi-square distribution with $d$ degrees of freedom, so that the density of $\|U\|$ is
\be\label{psi-normal}
\psi(r)=\frac{d^{d/2}}{2^{d/2-1}\,\Gamma(d/2)}\,r^{d-1}\,\e1^{-d r^2/2} \,, \ r \geq 0\,.
\ee
The distribution of $r$ is more and more concentrated around 1 as $d$ increases. Its $k$-th moment is
\bea
M_{\psi,k} = \Ex_\mu\{\|U\|^k\}=(2^{k/2}/d^{k/2}) \Gamma((k+d)/2)/\Gamma(d/2)\,, \ k=1,2,\ldots
\eea
so that its mean is smaller than 1 for any $d$ and its variance is $1/(2\,d)-1/(8\, d^2)+ \SO(1/d^3)$, $d\to\infty$.

When $\PP=\PP_a$ uniform on $\SS_{d-1}(a)$, the expected $(\mu,s)$-distortion $D_{\mu,s}(\PP_a^{[n]})$ is given by \eqref{DS-Pa},
with $\psi(\cdot)$ given by \eqref{psi-normal}.

\vsp
When $\PP=\mu_\ms$ is the probability measure of $\SN(\0b_d,\ms^2\Ib_d/d)$, $R=\|X\|$ has the p.d.f.\ $\phi_\ms(\rho)=(1/\ms)\psi(\rho/\ms)$ where $\psi(\cdot)$ given by \eqref{psi-normal}, and the expected $(\mu,s)$-distortion $D_{\mu,s}(\mu^{[n]}_\ms)$ is given by \eqref{Ls-mean-final}
where
\bea
H(r,t;\Phi,n) =\bigg(1-\frac{1}{\ms}\,\int_{\rho\geq 0} I_\upsilon(\delta,\delta) \,\psi(\rho/\ms)\,\dd\rho \bigg)^n \,,
\eea
with $\upsilon=\upsilon(t,\rho,r)$ given by \eqref{v}.

\vsp
Let $\ms^* = \ms^*(d,n,s)$ denote the value of $\ms$ that minimises the distortion $D_{\mu,s}(\mu^{[n]}_\ms)$ and $a^* = a^*(d,n,s)$ denote the value of $a$ that minimises $D_{\mu,s}(\PP_a^{[n]})$, as in previous sections.

The empirical distribution of an optimal quantiser $\Xb_n^*$ minimising $E_{\mu,s}(\Xb_n)$ weakly converges to $\SN(\0b_d,\ms_\infty^*\,\Ib_d/d)$ as $n\to\infty$, with $\ms_\infty^*=1+s/d$; see \cite[Table~7.1]{GrafL2000}. Figure~\ref{F:normal-a-sigma-s} illustrates that $\ms^*=\ms^*(d,n,s)$ increases with $s$, similarly to $\ms_\infty^*$, and that $a^*=a^*(d,n,s)$ also increases with $s$, contrary to Sections~\ref{S:sphere} and \ref{S:ball}. The reason is that large values of $r=\|U\|$ get increasing importance as $s$ increases, while $r=1$ in Section~\ref{S:sphere} and $r\leq 1$ in Section~\ref{S:ball}. Note that $\ms^*(d,n,s)$ is significantly smaller than $\ms_\infty^*=1+s/d$ (Figure~\ref{F:normal-a-sigma-s} is for $d=10$).

\begin{figure}[ht!]
\begin{center}
\includegraphics[width=.49\linewidth]{Rstar-sigmastar_mu-normal_n1000-d10}
\includegraphics[width=.49\linewidth]{Rstar-sigmastar_mu-normal_n100000-d10}
\end{center}
\caption{\small Optimal values $a^*(d,n,s)$ ({\color{Cerulean} $\bullet$}) and $\ms^*(d,n,s)$ ({\color{Dandelion} $\blacklozenge$}) as functions of $s$ for $d=10$, $n=1\,000$ (left) and $n=100\,000$ (right).}
\label{F:normal-a-sigma-s}
\end{figure}
% tst_quantization_random_normal.m

%Since $\mu_\ms$ is sub-Gaussian (see, e.g., \cite[Chap.~2,3]{Vershynin2018}), it satisfies \eqref{norm-concentration} and Proposition~\ref{Pr:norm-concentration} implies that $D_{\mu,s}(\mu_\ms^{[n]}) \to D_{\mu,s}(\PP_\ms^{[n]})$ when $d\to \infty$ for given $n$ and $s$.

Figure~\ref{F:normal} presents $D_{\mu,s}^{1/s}(\mu^{[n]}_\ms)$ and $D_{\mu,s}^{1/s}(\PP_{a=\ms}^{[n]})$ as functions of $\ms$ for $s=2$, $n=1\,000$ and $d=5,10$ and $20$.
Note that, already for $d=5$ (left panel), $1\,000$ points are not enough for $\ms^*=\ms^*(d,n,s)\simeq 1.1832$ to approach the asymptotically optimal value $\ms^*_\infty=1.4$. For $d\geq 10$ (center and right panels) the plots of $D_{\mu,s}^{1/s}(\PP_\ms^{[n]})$ and $D_{\mu,s}^{1/s}(\mu^{[n]}_\ms)$ are close, especially near their minima.
%for $s=2$ and $d=5$. For $d\geq 10$ (center and right panels), $n=1\,000$ points are not enough to approach the asymptotically optimal distribution and $D_{\mu,2}^{1/2}(\PP_{a^*}^{[n]})< D_{\mu,2}^{1/2}(\mu^{[n]}_{\ms^*})$, where $a^*=a^*(d,n,s)$ minimises $D_{\mu,s}(\PP_a^{[n]})$.
%As $d$ grows, $\psi(\cdot)$ is more and more concentrated around 1, and $\ms^*-a^*$ tends to zero: compare the three panels of Figure~\ref{F:normal}. For any $s\geq 1$ and large $d$, unless $n^{1/d}$ is large enough, $\PP_{a^*}$ provides a nearly optimal solution.
The behaviour is similar when using other values of $s$.


The left panel of Figure~\ref{F:normal2} shows that $\ms^*-a^*$ tends to zero as $d$ increases, a consequence of $D_{\mu,s}(\mu_\ms^{[n]})$ tending to $D_{\mu,s}(\PP_\ms^{[n]})$. Moreover, as $d$ grows, $\psi(\cdot)$ is more and more concentrated around 1. This explains why for $d\geq 10$ the values of $a^*$ are close to those obtained for $\mu$ uniform on $\SS_{d-1}(1)$; see the curves on Figure~\ref{F:sphere}-top-left (they are plotted for $s=1$ but those obtained for $s=2$ are very similar).
The right panel presents the efficiencies of the asymptotically optimal distribution $\mu_{\ms^*_\infty}$ relative to $\PP_{a^*}$ and $\mu_{\ms^*}$, for $n=10\,000$, as functions of $d$. These efficiencies are always less than 1, meaning that the asymptotic regime is not yet reached for the values of $d$ considered. At $d=10$, however, we can observe that, contrary to the central panel of Figure~\ref{F:normal} where $n=1\,000$, now $D_{\mu,2}(\PP_{a^*}^{[n]}) > D_{\mu,2}(\mu^{[n]}_{\ms^*})$: with $n=10\,000$ points we approach the asymptotic regime and $\PP_{a^*}$ is dominated by a suitably chosen normal distribution.

\begin{figure}[ht!]
\begin{center}
\includegraphics[width=.32\linewidth]{Es-normal-exact_R-sigma_n1000-d5-s2}
\includegraphics[width=.32\linewidth]{Es-normal-exact_R-sigma_n1000-d10-s2}
\includegraphics[width=.32\linewidth]{Es-normal-exact_R-sigma_n1000-d20-s2}
\end{center}
\caption{\small $D_{\mu,s}^{1/s}(\PP^{[n]}_{a=\ms})$ ({\color{Cerulean} {\bf - - -}}) and $D_{\mu,s}^{1/s}(\mu^{[n]}_\ms)$ ({\color{Dandelion} {\bf ---}}) as functions of $\sigma$ for different $d$ with $n=1\,000$ and $s=2$.}
\label{F:normal}
\end{figure}
% tst_quantization_random_normal.m

%\begin{comment}
%\begin{figure}[ht!]
%\begin{center}
%\includegraphics[width=.32\linewidth]{Es-normal-exact_R-sigma_n100000-d5-s4}
%\includegraphics[width=.32\linewidth]{Es-normal-exact_R-sigma_n100000-d10-s4}
%\includegraphics[width=.32\linewidth]{Es-normal-exact_R-sigma_n100000-d20-s4}
%\end{center}
%\caption{\small Same as Figure~\ref{F:normal} for $n=100\,000$ and $s=4$.}
%\label{F:normal_b}
%\end{figure}
%% tst_quantization_random_normal.m
%\end{comment}

\begin{figure}[ht!]
\begin{center}
\includegraphics[width=.49\linewidth]{Rstar-sigmastar_mu-normal_n10000-d10-40}
\includegraphics[width=.49\linewidth]{EffRstar-Effsigmastar_mu-normal_n10000-d10-40}
\end{center}
\caption{\small Left: Optimal values $a^*(d,n,s)$ ({\color{Cerulean} {\bf - - -}}) and $\ms^*(d,n,s)$ ({\color{Dandelion} {\bf ---}}) as functions of $d$.
Right: ratios $D_{\mu,s}^{1/s}(\PP^{[n]}_{a^*})/D_{\mu,s}^{1/s}(\mu^{[n]}_{\ms^*_\infty})$ ({\color{Cerulean} {\bf - - -}}) and $D_{\mu,s}^{1/s}(\mu^{[n]}_{\ms^*})/D_{\mu,s}^{1/s}(\mu^{[n]}_{\ms^*_\infty})$ ({\color{Dandelion} {\bf ---}}) as functions of $d$ for $s=2$ and $n=10\,000$.}
\label{F:normal2}
\end{figure}
% tst_quantization_random_normal.m

%\begin{comment}
%\begin{figure}[ht!]
%\begin{center}
%\includegraphics[width=.49\linewidth]{Rstar-sigmastar_mu-normal_n100000_s4_d10-40}
%\includegraphics[width=.49\linewidth]{EffRstar-Effsigmastar_mu-normal_n100000_s4_d10-40}
%\end{center}
%\caption{\small Same as Figure~\ref{F:normal2} for $n=100\,000$ and $s=4$.}
%\label{F:normal2_b}
%\end{figure}
%% tst_quantization_random_normal.m
%\end{comment}

%The same is true for every fixed $d\geq 3$: as $\mu_{\ms_\infty^*}$ is asymptotically optimal (for $n\to\infty$), $D_{\mu,s}^{1/s}(\PP_{a^*}^{[n]})$ is larger than $D_{\mu,s}^{1/s}(\mu^{[n]}_{\ms^*})$ when $n$ is large enough. However, since $D_{\mu,s}(\mu_\ms^{[n]}) \to D_{\mu,s}(\PP_\ms^{[n]})$ when $d\to \infty$, one may expect to have $D_{\mu,s}(\PP_{a^*}^{[n]}) < D_{\mu,s}(\mu_{\ms^*}^{[n]})$ for a given $n$ and large enough $d$. Proposition~\ref{Pr:norm-concentration} is for $n$ fixed, but the property $D_{\mu,s}(\mu_\ms^{[n]}) \to D_{\mu,s}(\PP_\ms^{[n]})$ remains valid when $n=n(d)$ grows slowly enough with $d$. Inspection of the proof of the theorem indicates that the condition is $n(d)\,\exp(- d\, \me^2) \to 0$ as $d\to\infty$ for any $\me>0$; that is, $n^{1/d}\to 1$. In fact, the numerical results presented below indicate that $D_{\mu,s}(\PP_{a^*}^{[n]}) < D_{\mu,s}(\mu_{\ms^*}^{[n]})$ even for $n$ growing exponentially fast with $d$.

In view of the right panel of Figure~\ref{F:normal2}, and following the analysis in Section~\ref{S:ball}, we investigate the following question: for a fixed dimension $d$, up to which sample size $n$ does the random quantiser $\PP_{a^*}^{[n]}$ outperform $\mu_{\ms^*}^{[n]}$? Defining
\be\label{nstar}
n^*(d,s) = \max\{n\in\NN: D_{\mu,s}(\PP_{a^*}^{[n]}) < D_{\mu,s}(\mu^{[n]}_{\ms^*}) \} \,,
\ee
in the right panel of Figure~\ref{F:nstar-ball_s2-4-10}
%Figure~\ref{F:nstar-normal_s2-4-10}
we present $n^*(d,s)$ (in log scale) as a function of $d=3,\ldots,20$ for three values of $s$. The straight lines fitted to these logarithmic plots show almost perfect coincidence between $n^*(d,s)$ and the exponential increase $n_0\, \ml^d$, with numerical values of $n_0$ and $\ml$ depending weakly on $s$: $n_0 \simeq 2.81$, 2.68, 2.21 and $\ml \simeq 2.08$, 2.07, 2.05 for $s=2$, 4 and 10, respectively.

%\begin{figure}[ht!]
%\begin{center}
%\includegraphics[width=.59\linewidth]{nstar-normal_d3-20_s2-4-10}
%\end{center}
%\caption{\small $n^*(d,s)$ given by \eqref{nstar} as a function of $d$ for $s=2$ ({\color{Dandelion} $\blacklozenge$}), $s=4$ ({\color{red} $\bigstar$}) and $s=10$ ({\color{Cerulean} $\bullet$}).}
%\label{F:nstar-normal_s2-4-10}
%\end{figure}
% tst_quantization_random_normal.m



%\begin{comment}
%
%It is remarkable that the performances of $\PP_{a^*}$ and $\PP_{M_{\psi,1}\widetilde{a^*}}$ are hardly distinguishable, as $\widetilde{a^*}$ is significantly easier to compute than $a^*$ ($D_{\mu,s}(\PP_a^{[n]})$ involves a single integration when $\mu$ is uniform on a sphere, see \eqref{Ds-sphere}, whereas a double integration is required when $\mu$ is normal, see \eqref{DS-Pa}). Assuming that $\widetilde\mu$ is uniform on $\SS_{d-1}(M_{\psi,1})$ rather than on $\SS_{d-1}(1)$ is also important for the asymptotic designs of Section~\ref{S:extreme-value-sphere}: compare the efficiencies obtained for $\PP_{\widehat{a^*_\ml}}$ (purple stars) and $\PP_{M_{\psi,1}\widehat{a^*_\ml}}$ (red squares).
%
%\begin{figure}[ht!]
%\begin{center}
%\includegraphics[width=.49\linewidth]{astar-sigstar-normal_n2d_d3-20_s2b}
%%\includegraphics[width=.49\linewidth]{Esmu-sigstar-normal_n2d_d3-20_s2}
%\includegraphics[width=.49\linewidth]{Eff-sigstar-normal_n2d_d3-20_s2b}
%\end{center}
%\caption{\small Left: $\ms^*(d,n,s)$ ({\color{Dandelion} $\blacklozenge$}), $M_{\psi,1}\widetilde{a^*}(d,n,s)\simeq a^*(d,n,s)$ ({\color{Cerulean} $\bullet$}), $\widetilde{a^*}(d,n,s)$ ({\color{green} $\blacktriangledown$}) and $M_{\psi,1}\widehat{a^*_\ml}$ ({\color{red} {\tiny $\blacksquare$}}) as functions of $d$; the value $\widehat{a^*_\ml}=\sqrt{3}/2$ is indicated by an horizontal dashed line;
%Right: efficiencies $D_{\mu,s}^{1/s}(\PP_a^{[n]})/D_{\mu,s}^{1/s}(\mu_{\ms^*}^{[n]})$ for $a=a^*(d,n,s)\simeq M_{\psi,1}\widetilde{a^*}(d,n,s)$ ({\color{Cerulean} $\bullet$}), $a=\widetilde{a^*}(d,n,s)$ ({\color{green} $\blacktriangledown$}), $a=\widehat{a^*_\ml}$ ({\color{Orchid} $\bigstar$}) and $a=M_{\psi,1}\widehat{a^*_\ml}$ ({\color{red} {\tiny $\blacksquare$}})
%%$D_{\mu,s}^{1/s}(\mu_{\ms^*}^{[n]})$ ({\color{Dandelion} $\blacklozenge$}), $D_{\mu,s}^{1/s}(\PP_{a^*}^{[n]})$ ({\color{Cerulean} $\bullet$}), $D_{\mu,s}^{1/s}(\PP_{\widetilde{a^*}}^{[n]})$ ({\color{green} $\blacktriangledown$}) and $D_{\mu,s}^{1/s}(\PP_{\widehat{a^*_\ml}}^{[n]})$ ({\color{Orchid} $\bigstar$})
%as functions of $d$ ($n=2^d$, $s=2$)
%{\Luc Check colors! Maybe Figure~\ref{F:normal_n2d_d3-20_s4} would be enough.}}
%\label{F:normal_n2d_d3-20_s2}
%\end{figure}
%% tst_quantization_random_normal.m
%
%Figure~\ref{F:normal_n2d_d3-20_s4} presents the same information as Figure~\ref{F:normal_n2d_d3-20_s2} but for $s=4$ instead of $s=2$. As $s$ increases, ZZZ
%
%\end{comment}


%\begin{comment}
%\begin{figure}[ht!]
%\begin{center}
%\includegraphics[width=.49\linewidth]{astar-sigstar-normal_n2d_d3-20_s10b}
%\includegraphics[width=.49\linewidth]{Eff-sigstar-normal_n2d_d3-20_s10b}
%\end{center}
%\caption{\small Same as in Figure~\ref{F:normal_n2d_d3-20_s4} but for $s=10$.}
%\label{F:normal_n2d_d3-20_s10}
%\end{figure}
%% tst_quantization_random_normal.m
%
%\end{comment}





%___________________________________________________________________________
%\section{quantisation: asymptotic results}\label{S:asymptotic}
\section{Extreme-value approximations}\label{S:asymptotic}


In this section, extreme-value theory is used to derive approximations of the $(\mu,s)$-distortion $D_{\mu,s}(\PP_a^{[n]})$, where $\PP_a$ is uniform on $\SS_{d-1}(a)$ (Section~\ref{S:Rn-uniform-on-SS}).
In Section~\ref{S:n-growing-with-d}, three asymptotic regimes are identified that govern the limiting behaviour of $D_{\mu,s}(\PP_a^{[n]})$ when $n$ increases with $d$.
The case where $\mu$ is uniform on $\SS_{d-1}(1)$ is considered first (Section~\ref{S:extreme-value-sphere}): we show that to each asymptotic regime is associated a limiting value of the optimal radius $a^*$ that does not depend on $s$. In Section~\ref{S:extreme-value-general-mu}, these results are extended to general spherically symmetric measures $\mu$ satisfying a norm-concentration property, with the uniform distribution in $\SB_d(1)$ and the multivariate normal distribution as illustrative examples.

%-----------------------------
%\subsection{Extreme-value approximation}\label{S:extreme-value}



%-----------------------------
\subsection{Uniform quantisers on a sphere}\label{S:Rn-uniform-on-SS}

Consider random quantisers $\Rb_n$ with distribution $\Px=\PP_a^{[n]}$ where $\PP_a$ is uniform on $\SS_{d-1}(a)$. For large $n$, the distribution of $\zeta(n,d)$ in \eqref{d2URn} can be approximated using standard results from extreme-value theory. Here we assume $d\geq 3$ to avoid singularity in beta distributions.

\begin{lemma}\label{L:main-extreme-value2}
Assume that $d\geq 3$ and $\Px=\PP_a^{[n]}$. For any fixed $\ub\in\RR^d$, we have
\be\label{d2URnb}
\frac{1}{4\,a\,\|\ub\|\,\kappa_{n,d}} \left[d^2(\ub,\Rb_n)-(\|\ub\|-a)^2\right] \rad \xi_d \,, \ n\to\infty\,,
\ee
where $\xi_d$ is a random variable with Weibull c.d.f.\ $F_d(t)=1-\exp(-t^\delta)$, $t\geq 0$, and moments $\Ex\{\zeta_d^k\}=\Gamma(1+k/\delta)$ with $\delta=(d-1)/2$, and where $\kappa_{n,d}$ is the $1/n$ quantile of $\beta_{\delta,\delta}$, defined by $I_{\kappa_{n,d}}(\delta,\delta)=1/n$.
\end{lemma}

\begin{proof}
From standard extreme-value theory (see, e.g., \cite[p.~59]{KotzN2000}), the random variable
$\zeta(n,d)/\kappa_{n,d}$ in \eqref{d2URn} converges in distribution to the random variable $\xi_d$.
\end{proof}

\begin{remark}\label{R:moments-duRn}
Lemma~\ref{L:main-extreme-value2} yields the following approximations for the $s$-moment of $d(\ub,\Rb_n)$ for large $n$ and any given $\ub$:
\bea
M_s(d,n,a,r) = \int_{\zeta\geq 0} \left[(r-a)^2 + 4\,a\,r\,\kappa_{n,d}\,\zeta\right]^{s/2}\, \dd F_d(\zeta) \,,
\eea
with $r=\|\ub\|$.
Explicit expressions are easily obtained for even $s$. Assuming that $\|\ub\|=1$, this gives in particular the approximations
\bea
\var_2 &=& 16\,a^2\,\kappa_{n,d}^2\,\var\{\xi_d\} = 16\,a^2\,\kappa_{n,d}^2 \left[\Gamma(1+2/\delta)-\Gamma^2(1+1/\delta)\right] \,, \\
\var_4 &=& 64\,a^2\,\kappa_{n,d}^2\, \left\{(1-a)^4\,\var\{\xi_d\} + a^2\,\kappa_{n,d}^2\,\var\{\xi_d^2\} \right. \\
&& \left. + (1-a)^2\,a\,\kappa_{n,d}\,\left[\Gamma(1+3/\delta)-\Gamma(1+1/\delta)\Gamma(1+2/\delta)\right] \right\} \,.
\eea
for $\var_s=\var\{d^s(\ub,\Rb_n)\}$ when $s=2$ and $s=4$, respectively.
\fin
\end{remark}

The random variables $\zeta(n,d)$ in \eqref{d2URn} depend on $\ub$ and $\Rb_n$, but \eqref{d2URnb} shows that the dependence in $\Rb_n$ vanishes asymptotically, so that $\xi_d$ only depends on $\ub/\|\ub\|$. We hence obtain the following extreme-value approximation of $D_{\mu,s}(\PP_a^{[n]})$ for $\mu$ spherically symmetric.

\begin{definition}\label{D:main-extreme-value}
Assume that $d\geq 3$ and that $\mu$ is spherically symmetric. For any $\Rb_n\simd \PP_a^{[n]}$, the extreme-value approximation of the $(\mu,s)$-distortion $E_{\mu,s}^s(\Rb_n)$ is defined by
\be\label{Ds-approx}
\widehat E_{\mu,s}^s(n;a) = \int_{r\geq 0} \int_{\zeta\geq 0} \left[(r-a)^2 + 4\,a\,r\,\kappa_{n,d}\,\zeta\right]^{s/2}\, \dd F_d(\zeta) \, \dd \Psi(r) \,,
\ee
where $\Psi(\cdot)$ is the c.d.f.\ of $\|U\|$ with $U\simd \mu$, $F_d(t)=1-\exp(-t^\delta)$ with $\delta=(d-1)/2$, and $\kappa_{n,d}$ is the $1/n$ quantile of $\beta_{\delta,\delta}$ defined by $I_{\kappa_{n,d}}(\delta,\delta)=1/n$. The extreme-value approximation $\widehat D_{\mu,s}(\PP_a^{[n]})$ of $D_{\mu,s}(\PP_a^{[n]})$ is defined by $\widehat D_{\mu,s}(\PP_a^{[n]})=\widehat E_{\mu,s}^s(n;a)$.
\end{definition}

The approximation \eqref{Ds-approx} is derived as follows. Since $\mu$ is spherically symmetric, we have $U\simd r\cdot U^{(d)}$ where $U^{(d)}$ is uniformly distributed on $\SS_{d-1}(1)$ and $r$ has the c.d.f.\ $\Psi(\cdot)$. Since $\xi_d$ in \eqref{d2URnb} only depends on $U^{(d)}$ and has the c.d.f.\ $F_d(\cdot)$, the substitution of $(r-a)^2 + 4\,a\,r\,\kappa_{n,d}\,\zeta(U^{(d)})$ for $d^2(U,\Rb_n)$ in $E_{\mu,s}^s(\Rb_n)= \Ex_\mu \{d^s(U,\Rb_n)\}$ gives \eqref{Ds-approx}. In contrast with Remark~\ref{R:moments-duRn} where $\ub$ is fixed, averaging over $\ub$ removes the variability (but the means are identical; that is, $\widehat E_{\mu,s}^s(n;a)=\Ex_{\PP_a^{[n]}}\{d^s(\ub,\Rb_n)\}$).

For $s$ even, $\widehat E_{\mu,s}^s(n;a)$ can be expressed explicitly in terms of moments of $\|U\|$, $M_{\Psi,k}=\Ex_\mu\{\|U\|^k\}=\int_{r\geq 0}  r^k\,\dd \Psi(r)$, and moments $\Ex\{\zeta^k\}=\Gamma(1+k/\delta)$ of the Weibull distribution. In particular, for $s=2$ and 4 we get the following extreme-value approximations for the $(\mu,s)$-distortions:
\be
\widehat E_{\mu,2}^2(n;a) &=& M_{\Psi,2} - 2\,a\,M_{\Psi,1}+a^2 + 4\,a\,\kappa_{n,d}\,M_{\Psi,1}\,\Gamma(1+1/\delta)  \,, \label{E2^2-general}\\
\widehat E_{\mu,4}^4(n;a) &=& M_{\Psi,4} - 4\,a\,M_{\Psi,3}+6\,a^2\,M_{\Psi,2}-4\,a^3\,M_{\Psi,1}+a^4 \nonumber \\
&& + \, 8\,a\,\kappa_{n,d}\,\Gamma(1+1/\delta)
\left( M_{\Psi,3} - 2\,a\,M_{\Psi,2}+a^2\,M_{\Psi,1} \right) \nonumber \\
&& + \, 16\,a^2\,\kappa_{n,d}^2\,M_{\Psi,2}\,\Gamma(1+2/\delta) \,. \label{E4^2-general}
\ee

Denote by $\widehat {a^*}=\widehat {a^*}(d,n,s)$ the value of $a$ that minimises $\widehat E_{\mu,s}(n;a)$, which forms an extreme-value approximation of the optimal $a^*$ minimising $D_{\mu,s}(\PP_a^{[n]})$. In the next corollary, we give the expressions of $\widehat {a^*}(d,n,2)$ and $\widehat E_{\mu,2}^2(n;\widehat {a^*})$; the later providing an extreme-value approximation for $\min_a D_{\mu,2}(\PP_a^{[n]})$.

\begin{corollary}\label{Coro:astar-s2-Pa}
For $d\geq 3$, $\mu$ spherically symmetric and $\Rb_n\simd \PP_a^{[n]}$, we have
\be
\widehat {a^*} = \widehat {a^*}(d,n,2) &=& M_{\Psi,1}\,\left[ 1 - 2\,\kappa_{n,d}\,\Gamma(1+1/\delta) \right]\,, \label{hat_a_2} \\
\widehat E_{\mu,2}^2(n;\widehat {a^*}) &=& M_{\Psi,2} - (\widehat {a^*})^2 \,. \nonumber
\ee
\end{corollary}

For $\mu$ uniform on $\SB_d(1)$, Figure~\ref{F:bstar_ALLastar_ball_s2_D} indicates that $\widehat {a^*}(d,n,2)$ and $a^*(d,n,2)$ are practically indistinguishable for $n=1\,000$ and $n=100\,000$ with $d=10,\ldots,50$.
For the situation considered in Section~\ref{S:normal} where $U\simd\SN(\0b_d,\Ib_d/d)$, $\widehat {a^*}(d,n,2)$ is also very close to $a^*(d,n,2)$ when $n$ is large enough; its evolution as a function of $d$ for $n=10\,000$ is visually indistinguishable from that of $a^*(d,n,s)$ on the left panel of Figure~\ref{F:normal2}.

The expressions for $\widehat {a^*} =\widehat {a^*}(d,n,4)$ and $\widehat E_{\mu,4}^4(n;\widehat {a^*})$ can also be written in analytic form (by finding the roots of a third-degree polynomial), but the expressions are cumbersome.

%--------------------------------------
\subsection{Three asymptotic regimes for $n$ growing with $d$}\label{S:n-growing-with-d}

In view of Lemma~\ref{L:main-extreme-value2}, the asymptotic behaviour of $\kappa_{n,d}$ is the key element in understanding that of $d^2(\ub,\Rb_n)$. For $d=3$, $\beta_{1,1}$ is uniform on $[0,1]$ and we simply have $\kappa_{n,3}=1/n$. For $d>3$ there is no explicit formula for $\kappa_{n,d}$, but it can easily be computed numerically by solving $I_{\kappa_{n,d}}(\delta,\delta)=1/n$. For fixed $t\leq 1/2$, $I_t(\delta,\delta)$ decreases as $\delta$ increases. Therefore, $\kappa_{n,d}$ is an increasing function of $d$ for any fixed $n\geq 2$ (it is also a decreasing function of $n$ for any fixed $d$). When $n$ tends to infinity with $d$ we have the following properties.

%\footnote{\AZ{Make this proposition a theorem?}. \Luc{We shall see\ldots It depends whether we can get strong results about other distributions than $\PP_a$.}}



%\begin{proposition} \label{th:kappa_n}
%For fixed $d$, the $1/n$ quantile $\kappa_{n,d}$ of  $\beta_{\delta,\delta}$ with $\delta=(d-1)/2$ tends to zero as $n$ tends to infinity. When $n=n(d)$ grows with $d\to\infty$, we have the following three cases:\\
%\noindent{\textit(i)} if $n^{1/d}\to\infty$, then $\lim_{d\to\infty} \kappa_{n,d} = 0$;\\
%\noindent{\textit(ii)} if $n=\ml^d\,(1+o(1))$ with $\ml>1$, then
%\be\label{bounds-kappa-caseii}
%\frac{1}{4\,\ml^2} \leq \liminf_{d\to\infty} \kappa_{n,d} \leq \limsup_{d\to\infty} \kappa_{n,d} \leq \frac12\,\left[1-\sqrt{1-1/\ml^2}\right]\,;
%\ee
%\noindent{\textit(iii)} if $n=d^\mt\,(1+o(1))$ with $\mt>0$, then $\lim_{d\to\infty} \kappa_{n,d} = \frac12$.
%\end{proposition}
%
%\begin{proof}
%Since $x^{\delta-1}(1-t)^\delta\leq x^{\delta-1}(1-t)^{\delta-1}\leq x^{\delta-1}(1-x)^{\delta-1}\leq x^{\delta-1}$ for $\delta\geq 1$ and all $t\in[0,1)$ and $x\in[0,t]$, we have
%\bea
%\frac{t^\delta (1-t)^\delta}{\delta B(\delta,\delta)} \leq I_t(\delta,\delta) \leq \frac{t^\delta}{\delta B(\delta,\delta)}\,.
%\eea
%This implies
%\be\label{bounds-kappa}
%c_d\,n^{-1/\delta} \leq \kappa_{n,d} \leq C_{n,d} = \left\{ \begin{array}{ll}
%\frac12 \left[1-\sqrt{1-4\,c_d\,n^{-1/\delta} }\right] & \mbox{ if } n > (4\,c_d)^\delta\,, \\
%\frac12 & \mbox{ otherwise}\,, \end{array}\right.
%\ee
%where $c_d=[\delta\,B(\delta,\delta)]^{1/\delta} = 1/4+ (\log d)/(4\,d)+\SO(1/d)$, $d\to\infty$, and $(4\,c_d)^\delta<d+2$ for all $d \geq 2$.
%
%When $d$ is fixed and $n\to\infty$, or $n^{1/d}\to\infty$ as $d\to\infty$, $n^{-1/d}\to 0$ implying that $C_{n,d}\to 0$. Therefore, $\kappa_{n,d}\to 0$.
%
%When $n=\ml^d\,(1+o(1))$ with $\ml>1$, by taking the limits on left and right hand sides of \eqref{bounds-kappa}, as $\lim_{d\to\infty} c_d=1/4$, we obtain \eqref{bounds-kappa-caseii}.
%
%Consider now the case $n=n(\mt,d)=d^\mt\,(1+o(1))$ with $\mt>0$. We show that the assumption $\liminf_{d\to\infty} \kappa_{n,d} = L < \frac12$ leads to a contradiction. Define $\ml_0=1/(2\,\sqrt{L})$, which is larger than 1. Take any $\ml\in(1,\ml_0)$. There exists $d_0$ such that for all $d>d_0$, $m=\lceil \ml^d \rceil > n(\mt,d)$, implying that $\kappa_{m,d}<\kappa_{n(\mt,d),d}$ and therefore $\liminf_{d\to\infty} \kappa_{m,d} \leq \liminf_{d\to\infty} \kappa_{n(\mt,d),d}=L$. However, \eqref{bounds-kappa-caseii} implies $\liminf_{d\to\infty} \kappa_{m,d}\geq 1/(4\,\ml^2)>L$, which gives a contradiction. Therefore, $\liminf_{d\to\infty} \kappa_{n,d} \geq 1/2$, and the right-hand side of \eqref{bounds-kappa} implies that $\lim_{d\to\infty} \kappa_{n,d} = 1/2$.
%\end{proof}

\begin{proposition} \label{th:kappa_n}
For fixed $d\geq 3$, the $1/n$ quantile $\kappa_{n,d}$ of  $\beta_{\delta,\delta}$ with $\delta=(d-1)/2$ tends to zero as $n$ tends to infinity. When $d \to \infty$ and $n=n(d)$ grows with $d$, we have the following three cases:\\
\noindent{\textit(i)} if $n^{1/d}\to\infty$, then $\lim_{d\to\infty} \kappa_{n,d} = 0$;\\
\noindent{\textit(ii)} if $n=C\ml^d\,(1+o(1))$ with $\ml>1$ and $C>0$, then
\be \label{asymptotic-kappa} \lim_{d\to\infty} \kappa_{n,d} = \frac12 \left[1-\sqrt{1-1/\ml^2}\right]\, ;
\ee
%\be\label{bounds-kappa-caseii}
%\frac{1}{4\,\ml^2} \leq \liminf_{d\to\infty} \kappa_{n,d} \leq \limsup_{d\to\infty} \kappa_{n,d} \leq \frac12\,\left[1-\sqrt{1-1/\ml^2}\right]\,;
%\ee
\noindent{\textit(iii)} %if $n=d^\mt\,(1+o(1))$ with $\mt>0$,
if $(\log n)/d \to 0$, then $\lim_{d\to\infty} \kappa_{n,d} = \frac12$.
\end{proposition}

\begin{proof}
Since $x^{\delta-1}(1-t)^\delta\leq x^{\delta-1}(1-t)^{\delta-1}\leq x^{\delta-1}(1-x)^{\delta-1}$ for $\delta\geq 1$ and all $t\in[0,1)$ and $x\in[0,t]$, we have
\be\label{It-ineq1}
\frac{t^\delta (1-t)^\delta}{\delta B(\delta,\delta)} \leq I_t(\delta,\delta) \,. %\leq \frac{t^\delta}{\delta B(\delta,\delta)}\,.
\ee
The derivation of an exploitable upper bound is more delicate.
Assume that $d \geq 3$ and let
$\delta'=\lfloor \delta \rfloor$, the largest integer smaller than or equal to $\delta$.
%Assume that $d\geq 3$ and let $\delta'$ be the largest integer $m$ such that $2\,m+1 \leq d$.
As $n\geq 2$, we are only interested in values of $t$ less than $1/2$. Since $\delta'\leq \delta$, we have $I_t(\delta,\delta) \leq I_t(\delta',\delta')$, which
%For any $n\geq 2$, we have $\kappa_{n,d}\leq \kappa_{n,2\,\delta'+1}$, where $\kappa_{n,2\,\delta'+1}$ is the solution of $I_t(\delta',\delta')=1/n$ with respect to $t$ and where the c.d.f.\ $I_t(\delta',\delta')$
can be calculated explicitly by successive integration by parts. Direct calculations yield the following special case of the well-known relation between the c.d.f.\ of the beta and binomial distributions:
\be\label{It1}
%I_t(\delta',\delta') = \frac{1}{\delta' B(\delta',\delta')} \, \sum_{j=1}^{\delta'} A_j\, t^{\delta'+j-1}(1-t)^{\delta'-j} \,,
I_t(\delta',\delta') &=& \sum_{k=0}^{\delta'-1} \binom{2\delta'-1}{k} t^{2\delta'-1-k}(1-t)^k \,, \\
&=& t^{\delta'}(1-t)^{\delta'-1} \sum_{k=0}^{\delta'-1} \binom{2\delta'-1}{k} \left(\frac{t}{1-t}\right)^{\delta'-1-k} \,.
\ee
%where
%\bea
%A_j= \frac{\delta'(\delta'-1)(\delta'-2)\cdots(\delta'-j+1)}{\delta'(\delta'+1)(\delta'+2)\cdots(\delta'+j-1)} = \frac{\delta'!(\delta'-1)!}{(\delta'-j)!(\delta'+j-1)!} \,.
%\eea
%We thus obtain
%\bea
%I_t(\delta',\delta') = \frac{t^{\delta'}(1-t)^{\delta'-1}}{\delta' %B(\delta',\delta')} \, \sum_{j=1}^{\delta'} A_j\, t^{j-1}(1-t)^{1-j} \leq \frac{t^{\delta'}(1-t)^{\delta'-1}}{\delta' B(\delta',\delta')} \sum_{j=1}^{\delta'} A_j \,,
%\eea
%where we used the property $t/(1-t) \leq 1$ for $t\leq 1/2$. Since $I_{1/2}(\delta',\delta')=1/2$, \eqref{It1} gives $\sum_{j=1}^{\delta'} A_j=4^{\delta'-1}\,\delta' B(\delta',\delta')$, so that
Using the property $t/(1-t) \leq 1$ for $t\leq 1/2$, with $\sum_{k=0}^{\delta'-1} \binom{2\delta'-1}{k}=4^{\delta'-1}$, we thus obtain
\be \label{upper_bound}
I_t(\delta,\delta) \leq I_t(\delta',\delta') \leq t^{\delta'}[4\,(1-t)]^{\delta'-1} \leq \frac12\,[4\,t(1-t)]^{\delta'-1}\,, \ t\leq 1/2\,.
\ee
%\footnote{\AZ{New: you simply removed one $t$ by using $t\leq 1$ but you could have used $t\leq 1/2$ and gained the multiplier $1/2$ in \eqref{upper_bound}. The lower bound in \eqref{bounds-kappa} would slightly improve. \\Also, to get the intermediate inequality you may perhaps use the inequality $1/(1-t)\leq 2 $ rather than $t/(1-t)\leq 1 $. In this way, I think we would  keep the correct value of the shape parameter in the intermediate inequality if $\delta=\delta'$. } }

Together with \eqref{It-ineq1}, this implies, for $n\geq 2$ and $d\geq 3$,
\be\label{bounds-kappa}
\frac12 \left[1-\sqrt{1- (2/n)^{1/(\delta'-1)}}\right]  \leq \kappa_{n,d} \leq \frac12 \left[1-\sqrt{(1-4\,c_d\,n^{-1/\delta})_+}\right] \,,
%c_d\,n^{-1/\delta} \leq \kappa_{n,d} \leq C_{n,d} = \left\{ \begin{array}{ll}
%\frac12 \left[1-\sqrt{1-4\,c_d\,n^{-1/\delta} }\right] & \mbox{ if } n > (4\,c_d)^\delta\,, \\
%\frac12 & \mbox{ otherwise}\,, \end{array}\right.
\ee
where $(x)_+=\max\{x,0\}$ and $c_d=[\delta\,B(\delta,\delta)]^{1/\delta} = 1/4+ (\log d)/(4\,d)+\SO(1/d)$, $d\to\infty$ (with, moreover, $(4\,c_d)^\delta<d+2$ for all $d \geq 2$, implying that the right-hand side of \eqref{bounds-kappa} is strictly smaller than  1/2 for $n\geq d+2$).

When $d$ is fixed and $n\to\infty$, or $n^{1/d}\to\infty$ as $d\to\infty$, $n^{-1/d}\to 0$ implying that the right-hand side of \eqref{bounds-kappa} tends to zero. Therefore, $\kappa_{n,d}\to 0$.

When $n=C\ml^d\,(1+o(1))$ with $\ml>1$ and $C>0$ (case (\textit{ii})), by taking the limits on left and right hand sides of \eqref{bounds-kappa}, as $\lim_{d\to\infty} c_d=1/4$, we obtain $\lim_{d\to\infty} \kappa_{n,d} = (1/2)\,\left[1-\sqrt{1-1/\ml^2}\right]$.

%When $n=n(\mt,d)=d^\mt\,(1+o(1))$ with $\mt>0$
When $\log n = o(d)$, (case (\textit{iii})), $n^{-1/d}\to 1$ and the left-hand side of \eqref{bounds-kappa} tends to 1/2.
\end{proof}

%\begin{comment}
%\AZ {
%\begin{remark}
%{\bf Alternative approach for getting an  upper bound like \eqref{upper_bound} with $\delta'=\delta$}
% By (26.5.14) in \cite{abramowitz1965handbook}, for $0<t<1/2$ we have
%$
%I_t(\delta,\delta) = \frac12 I_x(\delta,1/2)
%$, where $x=4t(1-t) <1$. By making the substitution $u=\sin^2 \theta $ in the integral $\int u^{\delta-1} (1-u)^{-1/2} \dd u$ (so that $\dd u= 2 \sin \theta \cos \theta $ and  $\sqrt{1-u^2}=\cos \theta$) we obtain
%\be \label{u_bound}
%I_t(\delta,\delta) = \frac12 I_x(\delta,1/2)= \frac {1}{B(\delta,1/2)} \int_0^{\arcsin (\sqrt{x})} \sin^{2 \delta -1} \theta \dd \theta  \, ,
%\ee
%where $B(\delta,1/2)=\sqrt{\pi} \Gamma (\delta)/\Gamma (\delta+1/2)$. For integer or half-integer $\delta$, this integral can be evaluated analytically; in particular, \eqref{It1} follows.\\
%To get an upper bound for $I_t(\delta,\delta)$, we can use, for example, the inequality $\sin \theta < \theta$ valid for all $\theta>0$. In this case,
%\eqref{u_bound} gives
%\bea \label{u_bound1}
%I_t(\delta,\delta) < \frac {1}{2 \delta B(\delta,1/2)} \left[\arcsin (\sqrt{4t(1-t)})\right] ^{2 \delta }
%= \frac {2^{2 \delta-1} \Gamma (\delta+1/2)}{ \sqrt{\pi} \, \delta\, \Gamma (\delta)} \arcsin^{2 \delta} (\sqrt{t})\, ,
%\eea
%where we have used the simplification $\arcsin (\sqrt{4t(1-t)})=2 \arcsin (\sqrt{t})$ valid for all $t \in (0,1/2)$.\\
%My new bound is very easy and has the right shape parameter of the extreme value distribution but unfortunately for me and this bound, it is worse than yours if $t$ is not close to 0. A much better upper bound is obtained  if we use the inequality $x-x^3/c\geq \sin x$ where $c=\pi^3/(\pi-2)\simeq 6.7901358$ but getting a closed-form expression for $\kappa_n$ is not possible.
%\end{remark}
%}
%\end{comment}
The left panel of Figure~\ref{F:sphere-approx0} plots $\kappa_{n,d}$ as a function of $d$ for different $n$: the continuation of each curve (with $n$ fixed) goes to the limiting value 1/2 in view of Proposition~\ref{th:kappa_n}-(\textit{iii}). The bounds in \eqref{bounds-kappa} are asymptotically accurate for large $n$, particularly the upper bound, but lack precision for smaller $n$. For illustration, only the envelope for $n = 10^3$ is shown.

%-----------------------------
\subsection{$\mu$ uniform on the unit sphere $\SS_{d-1}(1)$} \label{S:extreme-value-sphere}

When $U\simd \mu$ uniform on $\SS_{d-1}(1)$, for any $\Rb_n\simd \PP_a^{[n]}$ the approximations \eqref{E2^2-general} and \eqref{E4^2-general}
become
%%we have the approximation
%\be\label{approx-sphere-Es}
%\widehat E_{\mu,s}^s(n;a) = \widehat D_{\mu,s}(\PP^{[n]}_a) = \delta\,\int_{\zeta\geq 0} \left[(1-a)^2 + 4\,a\,\kappa_{n,d}\,\zeta\right]^{s/2}\, \zeta^{\delta-1} \e1^{-\zeta^\delta}\, \dd\zeta \,,
%\ee
%which can be calculated explicitly for any even $s$. For example, for $s=2$ and 4 we get the following extreme-value approximation for the quantisation errors:
\bea
\widehat E_{\mu,2}(n;a) \!\!\! &=& \!\!\! \left[ (1-a)^2 + 4\,a\,\kappa_{n,d}\,\Gamma(1+1/\delta) \right]^{1/2} \,, \\
\widehat E_{\mu,4}(n;a) \!\!\! &=& \!\!\! \left[ (1-a)^4 + 16\,a^2\,\kappa_{n,d}^2\,\Gamma(1+2/\delta) + 8\,a(1-a)^2\,\kappa_{n,d}\,\Gamma(1+1/\delta) \right]^{1/4} \,.
\eea
For $s=2$, Corollary~\ref{Coro:astar-s2-Pa} gives $\widehat {a^*}=\widehat {a^*}(d,n,2) = 1 - 2\,\kappa_{n,d}\,\Gamma(1+1/\delta)$ and $\widehat E_{\mu,2}^2(n;\widehat {a^*}) =1 - (\widehat {a^*})^2$.
The right panel of Figure~\ref{F:sphere-approx0} plots $\widehat {a^*}(d,n,2)-a^*(d,n,2)$ as a function of $d$ for different values of $n$: $\widehat {a^*}-a^* >0$ and the accuracy of the approximation of $a^*$ by $\widehat {a^*}$ tends to decrease with $d$ and increase with $n$.

\begin{figure}[ht!]
\begin{center}
\includegraphics[width=.49\linewidth]{Kappa_sphere_s2}
\includegraphics[width=.49\linewidth]{Rapprox-Rexact_sphere_s2}
\end{center}
\caption{\small Values of $\kappa_{n,d}$ (left) and of $\widehat {a^*}(d,n,2)-a^*(d,n,2)$ (right) as functions of $d$ for different values of $n$: $n=10^2$ ({\color{Dandelion} $\blacklozenge$}), $n=10^3$ ({\color{green} $\blacktriangledown$}), $n=10^4$ ({\color{Orchid} $\bigstar$}) and $n=10^5$ ({\color{gray} {\tiny $\blacksquare$}}).
}
\label{F:sphere-approx0}
\end{figure}
% tst_quantization_random_sphere.m


The left and right panels of Figure~\ref{F:sphere-approx} show that the two approximations $\widehat D_{\mu,2}(\PP^{[n]}_a)$ and $\widehat D_{\mu,4}(\PP^{[n]}_a)$ are quite accurate when $n$ is large enough for extreme value theory to be applicable. For large $d$, convergence to the extreme-value distribution is slow and $\widehat D_{\mu,s}^{1/s}(\PP^{[n]}_a)$ slightly overestimates the true value $D_{\mu,s}^{1/s}(\PP^{[n]}_a)$.

The left panel of Figure~\ref{F:sphere-approx-2} provides another illustration of the accuracy of the estimation, and presents normalised values $n^{1/d}\,\widehat D_{\mu,s}^{1/s}(\PP^{[n]}_a)$ and $n^{1/d}\,D_{\mu,s}^{1/s}(\PP^{[n]}_a)$ as functions of $n=2^m$, $m=6,7,\ldots,20$, for $s=4$ and $d=20$, both for $a=0.75$ fixed and $a$ optimal, i.e., $a^*(d,n,s)$ (the plots with $\widehat {a^*}(d,n,s)$ substituted for $a^*(d,n,s)$ are visually indistinguishable). The value $\widehat D_{\mu,s}^{1/s}(\PP^{[n]}_a)$ approaches $D_{\mu,s}^{1/s}(\PP^{[n]}_{a^*})$ for $n$ such that $a^*=a^*(d,n,s) \simeq 0.75$. On the right panel of Figure~\ref{F:sphere-approx-2}, for each $n$ considered, 100 values of $n^{1/d}\,E_{\mu,s}(\Rb_n)$ computed numerically\footnote{We used $2^{12}$ i.i.d.\ $U_i\simd \mu$, which provides an accurate enough estimation of $E_{\mu,s}(\Rb_n)$; see Remark~\ref{R:distortion-evaluation}.} for 100 random point sets $\Rb_n\simd \PP^{[n]}_{a^*}$ have been added to the plots of $n^{1/d}\,\widehat D_{\mu,s}^{1/s}(\PP^{[n]}_{a^*})$ and $n^{1/d}\,D_{\mu,s}^{1/s}(\PP^{[n]}_{a^*})$. The simulations are in perfect adequation with the exact mean value $n^{1/d}\,D_{\mu,s}^{1/s}(\PP^{[n]}_{a^*})$ and indicate that the variability of $E_{\mu,s}(\Rb_n)$ across quantisers is very small.


\begin{figure}[ht!]
\begin{center}
\includegraphics[width=.49\linewidth]{Esmu_sphere_s2_ralative_exact-approx}
\includegraphics[width=.49\linewidth]{Esmu_sphere_s4_ralative_exact-approx} \\
\end{center}
\caption{\small Relative error $1-\widehat D_{\mu,s}^{1/s}(\PP^{[n]}_a)/D_{\mu,s}^{1/s}(\PP^{[n]}_a)$ of the approximation of $D_{\mu,s}^{1/s}(\PP^{[n]}_a)$ (left for $s=2$ and right for $s=4$) based on extreme-value theory, as a function of $d$ for different values of $n$: $n=10^2$ ({\color{Dandelion} $\blacklozenge$}), $n=10^3$ ({\color{green} $\blacktriangledown$}), $n=10^4$ ({\color{Orchid} $\bigstar$}) and $n=10^5$ ({\color{gray} {\tiny $\blacksquare$}}); $a$ is the optimal value $a^*(d,n,s)$.
}
\label{F:sphere-approx}
\end{figure}
% tst_quantization_random_sphere.m



\begin{figure}[ht!]
\begin{center}
\includegraphics[width=.49\linewidth]{Es-true-approx-d20-s4-astar-a075}
\includegraphics[width=.49\linewidth]{Es-true-approx-d20-s4-astar-random}
\end{center}
\caption{\small $n^{1/d}\,\widehat D_{\mu,s}^{1/s}(\PP^{[n]}_a)$ ({\color{Dandelion} $\blacklozenge$}) and $n^{1/d}\,D_{\mu,s}^{1/s}(\PP^{[n]}_a)$ ({\color{Cerulean} $\bullet$}) as functions of $n$ for $s=4$ and $d=20$. Left: $a=0.75$ for top curves and $a=a^*(d,n,s)$ for bottom curves. Right: $a=a^*(d,n,s)$, the values of $n^{1/d}\,E_{\mu,s}(\Rb_n)$ computed numerically for 100 random quantisers $\Rb_n\simd \PP^{[n]}_{a^*}$ are shown as magenta dots.
}
\label{F:sphere-approx-2}
\end{figure}
% tst_quantization_random_sphere.m

\vsp
The property below shows that the asymptotic
behaviour of $\widehat {a^*}(d,n,s)$ when $n$ tends to infinity with $d$ does not depend on $s$.

\begin{corollary}\label{Coro:astar-asymptotic}
When $n=n(d)$ grows with $d\to\infty$, we have the following three cases for the limit of $\widehat {a^*}(d,n,s)$:\\
\noindent{\textit(i)} if $n^{1/d}\to\infty$, then $\lim_{d\to\infty} \widehat {a^*}(d,n,s) = 1$;\\
\noindent{\textit(ii)} if $n=\ml^d\,(1+o(1))$ with $\ml>1$, then $\lim_{d\to\infty} \widehat {a^*}(d,n,s) = \widehat {a^*_\ml} = \sqrt{1-1/\ml^2}$;\\
\noindent{\textit(iii)} if $\log n = o(d)$, then $\lim_{d\to\infty} \widehat {a^*}(d,n,s) = 0$.
\end{corollary}

\begin{proof}
Definition \eqref{Ds-approx} gives
\bea
\widehat E_{\mu,s}^s(n;a) = \int_{\zeta\geq 0} \left[(1-a)^2 + 4\,a\,\kappa_{n,d}\,\zeta\right]^{s/2}\, \dd F_d(\zeta) \,,
\eea
where the Weibull distribution tends to be concentrated at $1$ when $d\to\infty$. This implies that, for any given $s$, $\lim_{d\to\infty} \widehat E_{\mu,s}(n;a) - [(1-a)^2 + 4\,a\,\kappa_{n,d}]=0$, where $\kappa_{n,d}$ has the limiting behaviour indicated in Proposition~\ref{th:kappa_n}. Therefore,  $\lim_{d\to\infty} \widehat {a^*}(d,n,s) = \lim_{d\to\infty} 1-2\,\kappa_{n,d}$ does not depend on $s$ and is as indicated above for the three cases considered.
\end{proof}

On the left panel of Figure~\ref{F:sphere} we can observe the tendency $a^*(d,n,s)\to 0$ when $d$ is large but $n$ is not exponentially large compared to $d$; on the top-left panel with $d=10$ with see that $a^*(d,n,1)$ approaches 1 as $n$ increases. The left panel of Figure~\ref{F:sphere3} provides an illustration of  case (\textit{ii}) for $n=2^d$ (i.e., $\ml=2$) with $s=2,4$, for which $\widehat {a^*_\ml} = \sqrt{3}/2 \simeq 0.8660$.

%When $n^{1/d}$ is large, $\kappa_{n,d}$ is close to zero from Proposition~\ref{th:kappa_n}-(\textit{i}), and then $\widehat {a^*}(d,n,s)$ is close to 1 for any $s$.
%However, if $d$ is large but $n$ is not exponentially large compared to $d$, then Proposition~\ref{th:kappa_n}-(\textit{iii}) indicates that $\kappa_{n,d}$ is close to 1/2;
%the Weibull random variable $\xi_d$ of Lemma~\ref{L:main-extreme-value2} is almost concentrated at $1$, and $\widehat {a^*}(d,n,s)$ is close to zero for any $s$. This is the tendency we can observe on the left panel of Figure~\ref{F:sphere} for large $d$. Intermediate situations, when $n$ grows like $\ml^d$ with $\ml>1$, $\widehat {a^*}(d,n,s)$ tends to a limit $\widehat {a^*_\ml}$ that does not depend on $s$ and can be deduced from Proposition~\ref{th:kappa_n}-(\textit{ii}). The left panel of Figure~\ref{F:sphere3} provides an illustration for the case $n=2^d$, with $s=2,4$, for which $\widehat {a^*_\ml} = \sqrt{3}/2 \simeq 0.8660$.

\vsp
From Lemma~\ref{L:main-extreme-value2}, the same arguments as those leading to Definition~\ref{D:main-extreme-value} indicate that, for any $\Rb_n\simd \PP_a^{[n]}$ the extreme-value approximation of the $\mg$-quantile of the  distance c.d.f.\ $F_n(t; \Rb_n,\mu)$ is $\widehat q_\mg = \sqrt{(1-a)^2+4\,a\,\kappa_{n,d}\,t_\mg}$, with $t_\mg=[-\log(1-\mg)]^{1/\delta}$ the $\mg$-quantile of the Weibull distribution with c.d.f.\ $F_d(t)=1-\exp(-t^\delta)$, $\delta=(d-1)/2$.

The minimum of $\widehat q_\mg$ with respect to $a$ is obtained for $\widehat {a_*}(d,n,\mg)=1-2\,\kappa_{n,d}\,t_\mg$. Since $t_\mg\to 1$ as $d\to\infty$, for any fixed $\mg$, $\widehat {a_*}(d,n,\mg)$ has the same limits as those indicated in Corollary~\ref{Coro:astar-asymptotic} for $\widehat {a^*}(d,n,s)$, independently of $\mg$, when $n = n(d) \to \infty.$

\begin{remark}\label{R:quantiles} Note that $\widehat q_\mg = \widehat E_{\mu,2}(n;a)$ for any $a$ when $\mg=\mg(d,2)=F_d[\Gamma(1+1/\delta)]$, showing that for large $n$, minimising the expected $(\mu,2)$-distortion is equivalent to minimising the $\mg$-quantile of the mean distance c.d.f.\ for
\bea
\mg=\mg(d,2)=1- \e1^{-\e1^{-\mgbb}} + \frac{\pi^2}{6\,d}\, \e1^{-\mgbb-\e1^{-\mgbb}} +\SO(d^{-2})  \,, \ d\to\infty\,,
\eea
where $\mgbb$ is Euler's constant and $1- \e1^{-\e1^{-\mgbb}}\simeq 0.429624$. This suggests that $L_2$-quantisation is roughly equivalent to minimising the median of the distance c.d.f., a phenomenon that can also be observed for other sets than the sphere.
\fin
\end{remark}

%-----------------------------
\subsection{$\mu$ possesses a norm-concentration property} \label{S:extreme-value-general-mu}

We say that the probability measure  $\mu$ on $\RR^d$ possesses a {\em norm-concentration property} when it satisfies
\be\label{norm-concentration}
\mbox{for all } \me>0\,, \quad \mu\left\{\left| \|U\|-A \right| \geq \me\right\} \leq 2\,\exp(-c\, d\, \me^2)
\ee
when $U\sim\mu$, for some $A>0$ and $c$ a constant not depending on $d$. A multivariate distribution such that the components of $U$ are independent and sub-Gaussian with zero mean and strictly positive variance satisfies \eqref{norm-concentration}, see \cite[Th.~3.1.1]{Vershynin2018}. Here we consider distributions that satisfy \eqref{norm-concentration} and are spherically-symmetric, with the uniform distribution in $\SB_b(1)$ and the multivariate normal distribution for which $r=\|U\|$ has the density \eqref{psi-normal} as particular cases. Both satisfy \eqref{norm-concentration} with $A=1$ and, without any loss of generality we consider measures $\mu$ such that $A=1$ in the following.

We continue to consider random quantisers $\Rb_n\simd\PP_a^{[n]}$, focusing on the extreme-value approximation $\widehat D_{\mu,s}(\PP_a^{[n]})$ of the $(\mu,s)$-distortion as defined in Definition~\ref{Ds-approx}. When $n=n(d)$ grows with $d$ and $\mu$ is simply spherically symmetric, the limiting value of $\widehat {a^*}(d,n,s)$ which minimises $\widehat D_{\mu,s}(\PP_a^{[n]})$  generally depends on $s$ in cases (\textit{i}) and (\textit{ii}) of Corollary~\ref{Coro:astar-asymptotic}.  The next proposition shows that, if $\mu$ satisfies \eqref{norm-concentration}, the asymptotic behaviour of $\widehat{a^*}(d,n,s)$ remains unchanged from that described in Corollary~\ref{Coro:astar-asymptotic}.

\begin{proposition} \label{Prop:extreme-value-general-mu}
Assume that $d\geq 3$, that $\Px=\PP_a^{[n]}$, and that $\mu$ is spherically symmetric and satisfies \eqref{norm-concentration}. Then, as $d \to \infty$ with $n = n(d)$ growing accordingly, the limit of $\widehat{a^*}(d,n,s)$ follows the three cases outlined in Corollary~\ref{Coro:astar-asymptotic}.
\end{proposition}

\begin{proof} Consider the random variable $T_{n,d}=(r-a)^2 + 4\,a\,r\,\kappa_{n,d}\,\zeta$ in \eqref{Ds-approx}. The limiting value $\kappa_\infty$ of the quantile $\kappa_{n,d}$ in the three cases considered is given in Proposition~\ref{th:kappa_n}, and $T_{n,d}$ converges in probability to $T_\infty=(1-a)^2 + 4\,a\,\kappa_\infty$ when $d\to\infty$. From \eqref{norm-concentration}, for any given $s>0$ the sequence $\{T_{n,d}^{s/2}\}$ is bounded in $L^p$ for any $p$ and is thus uniformly integrable. Therefore, $\widehat D_{\mu,s}(\PP_a^{[n]})\to T_\infty^{s/2}$, which is minimum for the values of $a$ indicated in Corollary~\ref{Coro:astar-asymptotic}.
\end{proof}

%However, when $\mu$ is uniform on $\SB_d(1)$ or is spherically normal with $\|U\|$ having the density \eqref{psi-normal}, for any fixed $k$ the moment $\Ex\{\|U\|^k\}$ tends to 1 as $d\to\infty$, so that for any given $s$, $\lim_{d\to\infty} \widehat E_{\mu,s}(n;a) -[(1-a)^2 + 4\,a\,\kappa_{n,d}]=0$. The asymptotic behaviour of $\widehat {a^*}(d,n,s)$ when $n$ tends to infinity with $d$ is thus the same as in Corollary~\ref{Coro:astar-asymptotic}.

%The norm-concentration property has the consequence that for large $d$, $D_{\mu,s}(\PP^{[n]})$ is close to $D_{\mu,s}(\PP_\ms^{[n]})$ where $\PP_\ms$ is uniform on the sphere $\SS_{d-1}(\ms)$. This property, which follows from Proposition~\ref{th:dist} is stated below. We denote
%\bea
%J_{\ms,s} = \int_{r\geq 0} (|r-\ms|+2\,r)^{s/2}\,\dd\Psi(r) \,.
%\eea
%and $C_s=(s/2)\,\max\{2^{s/2-2},1\}$.

%\begin{comment}
%\begin{proposition}\label{Pr:norm-concentration}
%Let the distribution $\PP$ satisfy the norm-concentration property \eqref{norm-concentration}. Then for any $s\geq 2$, any $n$ such that $D_{\mu,s}(\PP_\ms^{[n]})\leq J_{\ms,s}^2$ and any $\me$ satisfying
%\bea
%\me\leq \mO_{\mu,s,\ms,n}=\frac12\,\left[\frac{D_{\mu,s}(\PP_\ms^{[n]})}{J_{\ms,s}}\right]^{2/s}\,,
%\eea
%we have
%\bea
%\left|\frac{D_{\mu,s}(\PP^{[n]})}{D_{\mu,s}(\PP_\ms^{[n]})}-1\right| \leq 2\,\me\, \frac{\max\{2^{s/2-2},1\}}{\mO_{\mu,s,\ms,n}}
%\eea
%with probability
%at least $P_{n,d}(\me)=1-2\,n\exp(-c\,d\,\me^2)$.
%\end{proposition}
%
%\begin{proof}
%ZZZ THE PROOF IS NOT VALID!
%
%
%Denote $\Delta_r(\rho,z)=(r-\rho)^2+4\,r\,\rho\,z$. Identity \eqref{eq:quantile2} of Proposition~\ref{th:dist} implies
%\bea
%D_{\mu,s}(\PP^{[n]}) = \int_{r\geq 0} \Ex\left\{ \min_{i=1,\ldots,n} \Delta_r^{s/2}(R_i,\zeta_i) \right\}\, \dd\Psi(r) \,,
%\eea
%where the expectation is with respect to the $R_i$ and $\zeta_i$ which are independent and have the c.d.f.\ $\Phi(\cdot)$ and the p.d.f.\ $\varphi_{\delta,\delta}(\cdot)$, respectively.
%
%Since $\PP$ satisfies \eqref{norm-concentration}, we have
%\bea
%\mbox{for all } \me>0\,, \quad \Pr\left\{ \max_{i=1,\ldots,n} \left| R_i-\ms \right| \geq \me\right\} \leq 2\,n\,\exp(-c\, d\, \me^2)\,,
%\eea
%so that, with probability $P(n,d) \geq 1-2\,n\,\exp(-c\, d\, \me^2)$, $\ms-\me< R_i< R_i+\me$ for all $i$. In that case, direct calculation shows that
%\bea
%\Delta_r(\ms,\zeta_i)-h_\me(\ms,r) \leq \Delta_r(R_i,\zeta_i) \leq \Delta_r(\ms,\zeta_i)+h_\me(\ms,r)\,, \ i=1,\ldots,n\,,
%\eea
%and therefore
%\bea
%\left|\min_{i=1,\ldots,n} \Delta_r(R_i,\zeta_i) - \min_{i=1,\ldots,n} \Delta_r(\ms,\zeta_i) \right| \leq h_\me(\ms,r)\,,
%\eea
%where $h_\me(\ms,r)=\me^2+2\,\me\,|r-\ms|+4\,\me\,r$. When $s\geq 2$, this implies
%\bea
%\left|D_{\mu,s}(\PP^{[n]}) - D_{\mu,s}(\PP_\ms^{[n]})\right| && \\
%&& \hspace{-3cm} \leq \frac{s}{2} \int_{r\geq 0} \Ex\left\{ \left[ \min_{i=1,\ldots,n} \Delta_r(\ms,\zeta_i) + h_\me(\ms,r)\right]^{s/2-1}\right\} h_\me(\ms,r)\, \dd\Psi(r)\,, \\
%&& \hspace{-3cm} \leq s\,2^{s/2-2} \int_{r\geq 0} \left[ D_{\PP_r,s/2-1}(\PP_\ms^{[n]})\, h_\me(\ms,r) + h_\me^{s/2}(\ms,r) \right] \dd\Psi(r)\,.
%\eea
%%where the expectation is with respect to the $\zeta_i$.
%The right-hand side tends to zero with $\me$ and $P(n,d)$ tends to 1 as $d\to\infty$, implying that $D_{\mu,s}(\PP^{[n]})$ tends to $D_{\mu,s}(\PP_\ms^{[n]})$ when $d\to \infty$.
%\end{proof}
%\end{comment}

To illustrate the above analysis, we consider again two specific cases for $\mu$: uniform distribution in the unit ball $\SB_d(1)$, in the small $n$ regime with fixed $n$ and increasing $d$; multivariate normal distribution $\SN(\0b_d, \Ib_d/d)$, in the exponential regime with $n = 2^d$.


%-----------------------------
\subsubsection{$\mu$ is uniform in the unit ball $\SB_d(1)$} \label{S:extreme-value-ball}

%Definition~\ref{D:main-extreme-value} gives an extreme-value approximation of the $(\mu,s)$-distortion $E_{\mu,s}^s(\Rb_n)$ of a random quantiser $\Rb_n$ distributed with $\PP_a^{[n]}$. This yields an approximation $\widehat{a^*}(d,n,s)$ for the optimal choice of $a$: for $s=2$,

In the non-asymptotic regime, $\widehat{a^*}(d,n,2)$ is given by \eqref{hat_a_2} for $s=2$; for $s=4$, $\widehat{a^*}(d,n,4)$ is a root of the third-degree polynomial corresponding to the derivative of \eqref{E4^2-general}.
%A second-level asymptotic approximation can be obtained by substituting $\lambda=n^{1/d}$ in the asymptotic expression \eqref{asymptotic-kappa} for $\kappa{n,d}$. For example, when $s=2$, this yields the approximation $\widehat{\widehat{a^*}}(d,n,2)=M_{\Psi,1} \sqrt{1-n^{-2/d}}\,\Gamma[1+2/(d-1)]$; see \eqref{hat_a_2}.

Figure~\ref{F:bstar_ALLastar_ball_s2_D} displays the following quantities as functions of $d$ for two values of $n$: $a^*(d,n,2)$, its approximation $\widehat{a^*}(d,n,2)$,
%its second-level approximation $\widehat{\widehat{a^*}}(d,n,2)$,
and $b^*(d,n,2)$, which minimises $D_{\mu,2}(\PP_{0,b}^{[n]})$ (see Section~\ref{S:ball}). The comparative behaviour of $a^*(d,n,2)$ and $b^*(d,n,2)$ has already been discussed in Section~\ref{S:ball}; the plots of $a^*(d,n,2)$ and $\widehat{a^*}(d,n,2)$ are practically indistinguishable.
%$\widehat{\widehat{a^*}}(d,n,2)$ differs significantly. However, Figure~\ref{F:ALLeff_ball_s2_D} demonstrates that the increase in distortion resulting from the use of the second-level approximation $\widehat{\widehat{a^*}}(d,n,2)$ instead of $\widehat{a^*}(d,n,2)$ (or the exact value $a^*(d,n,2)$) is very small.

Figure~\ref{F:ALLeff_ball_s2_D} displays the efficiencies $D_{\mu,s}^{1/s}(\PP_a^{[n]})/D_{\mu,s}^{1/s}(\PP_{0,b^*}^{[n]})$ as functions of $d$ for $a=a^*(d,n,2)$ and $a=\widehat{a^*}(d,n,2)$. These efficiencies are quasi indistinguishable, and $\PP_{a^*}^{[n]}$ and $\PP_{\widehat{a^*}}^{[n]}$ both outperform $\PP_{0,b^*}^{[n]}$ for the values of $n$ and $d$ considered. Additionally, $\PP_{0,b^*}^{[n]}$ significantly outperforms $\PP_{0,1}^{[n]}$; see Figure~\ref{F:ball} in Section~\ref{S:ball}.
%($\PP_{\widehat{\widehat{a^*}}}^{[n]}$ is slightly worse than $\PP_{0,b^*}^{[n]}$ but still much better than $\PP_{0,1}^{[n]}$).

\begin{figure}[ht!]
\begin{center}
\includegraphics[width=.49\linewidth]{bstar_ALLastar_ball_s2_n1000_D}
\includegraphics[width=.49\linewidth]{bstar_ALLastar_ball_s2_n100000_D}
\end{center}
\caption{\small $a^*(d,n,2)$ ({\color{Cerulean} $\bullet$}), $\widehat{a^*}(d,n,2)$ ({\color{red} $\bigstar$})
%, $\widehat{\widehat{a^*}}(d,n,2)$ ({\color{green} $\blacktriangledown$})
and $b^*(d,n,2)$ ({\color{Dandelion} $\blacklozenge$}) as functions of $d$ for $n=1\,000$ (left) and $n=100\,000$ (right).}
\label{F:bstar_ALLastar_ball_s2_D}
\end{figure}
% tst_quantization_random_ball.m


\begin{figure}[ht!]
\begin{center}
\includegraphics[width=.49\linewidth]{ALLeff_ball_s2_n1000_D}
\includegraphics[width=.49\linewidth]{ALLeff_ball_s2_n100000_D}
\end{center}
\caption{\small Efficiencies $D_{\mu,s}^{1/s}(\PP_a^{[n]})/D_{\mu,s}^{1/s}(\PP_{0,b^*}^{[n]})$ as functions of $d$: $a=a^*(d,n,2)$ ({\color{Cerulean} $\bullet$}), $a=\widehat{a^*}(d,n,2)$ ({\color{red} $\bigstar$});
%, $\widehat{\widehat{a^*}}(d,n,2)$ ({\color{green} $\blacktriangledown$})
$n=1\,000$ (left) and $n=100\,000$ (right).}
\label{F:ALLeff_ball_s2_D}
\end{figure}
% tst_quantization_random_ball.m

%\begin{remark} When $\mu$ is simply spherically symmetric, the limiting value of $\widehat {a^*}(d,n,s)$ in cases (\textit{i}) and (\textit{ii}) generally depends on $s$. However, when $\mu$ is uniform on $\SB_d(1)$ or is spherically normal with $\|U\|$ having the density \eqref{psi-normal}, for any fixed $k$ the moment $\Ex\{\|U\|^k\}$ tends to 1 as $d\to\infty$, so that for any given $s$, $\lim_{d\to\infty} \widehat E_{\mu,s}(n;a) -[(1-a)^2 + 4\,a\,\kappa_{n,d}]=0$. The asymptotic behaviour of $\widehat {a^*}(d,n,s)$ when $n$ tends to infinity with $d$ is thus the same as in Corollary~\ref{Coro:astar-asymptotic}.
%{\Luc I would like to say something connected with optimal transport (Kantorovitch distance) for Gaussians (the fact that for large $d$, if $n$ is not astronomically large the "best sample" should not be taken Gaussian).}
%\fin
%\end{remark}

%For any fixed $k$ the moment $M_{\Psi,k}=\Ex_\mu\{\|U\|^k\}=1-k/(d+k)$ tends to 1 as $d\to\infty$. Therefore, although the value of $\widehat {a^*}(d,n,s)$ depends on $s$, $\lim_{d\to\infty} \widehat E_{\mu,s}(n;a) -[(1-a)^2 + 4\,a\,\kappa_{n,d}]=0$ for any given $s$. The asymptotic behaviour of $\widehat {a^*}(d,n,s)$ when $n$ tends to infinity with $d$ is thus the same as indicated in Corollary~\ref{Coro:astar-asymptotic}.

%-----------------------------
\subsubsection{$\mu$ is normal $\SN(\0b_d,\Ib_d/d)$} \label{S:extreme-value-normal}

We noticed in Section~\ref{S:normal} that $D_{\mu,s}(\PP_{a^*}^{[n]}) < D_{\mu,s}(\mu^{[n]}_{\ms^*})$ for $n \lesssim n_0\, \ml_0^d$ for some $n_0$ and $\ml_0$, where $a^*=a^*(d,n,s)$ and $\ms^*$ respectively minimise $D_{\mu,s}(\PP_a^{[n]})$ and $D_{\mu,s}(\mu^{[n]}_{\ms})$. Here we investigate the behaviour of $D_{\mu,s}(\PP_{a}^{[n]})$ for other choices of $a$ in a similar asymptotic regime where $n$ grows exponentially fast with $d$.

Given that $M_{\psi,1}=\Ex\{\|U\|\} \to 1$ and $\var\{\|U\|\} \to 0$ as $d\to\infty$ when $U \sim \mu$, we analyse the following cases for $a$.
(\textit{i}) Uniform distribution on $\SS_{d-1}(1)$:
we consider $a = \widetilde{a^*}(d,n,s)$ which minimises $D_{\widetilde\mu,s}^{1/s}(\PP_a^{[n]})$ for $\widetilde\mu$ uniform on $\SS_{d-1}(1)$, as discussed in Section~\ref{S:sphere}.
(\textit{ii}) Scaled uniform distribution: we also consider $a = M_{\psi,1}\widetilde{a^*}(d,n,s)$, obtained when $\widetilde\mu$ is uniform on $\SS_{d-1}(M_{\psi,1})$.
As we consider the asymptotic regime $n = 2^d$, we additionally examine (\textit{iii}) the limiting value $\widehat{a^*_\ml} = \sqrt{3}/2$ from Corollary~\ref{Coro:astar-asymptotic}-(\textit{ii});
(\textit{iv}) the value $\widehat{a^*}(d,n,s)$, which minimises $\widehat{E}_{\mu,s}^s(n;a)$ given by \eqref{Ds-approx}, where the quantile $\kappa_{n,d}$ is replaced by its asymptotic value $(1/2)(1 - \sqrt{3}/2)$.


 Figure~\ref{F:normal_n2d_d3-20_s4} (left) displays $\ms^*(d,n,s)$, $a^*(d,n,s)$, $\widetilde{a^*}(d,n,s)$ and $M_{\psi,1}\widetilde{a^*}(d,n,s)$ as functions of $d$ for $s=4$, together with $\widehat{a^*_\ml}$ (indicated by an horizontal line) and
$\widehat{a^*}(d,n,s)$.
%$M_{\psi,1}\widehat{a^*_\ml}$
The right panel in the same figure shows the efficiency of the optimal normal distribution $\mu_{\ms^*}$ compared to $\PP_a$ uniform on $\SS_{d-1}(a)$, that is, $D_{\mu,s}^{1/s}(\PP_a^{[n]})/D_{\mu,s}^{1/s}(\mu_{\ms^*}^{[n]})$, for the five choices of $a$ considered: $\PP_{a}$ performs better than $\mu_{\ms^*}$ (for $d\geq 4$ when $a=M_{\psi,1}\widetilde{a^*}(d,n,s)$). Note that the efficiencies for $\widehat{a^*_\ml}$ (purple triangles) and $a^*(d,n,s)$ (blue dots) almost coincide for $d\geq 5$.

\begin{figure}[ht!]
\begin{center}
\includegraphics[width=.49\linewidth]{astar-sigstar-normal_n2d_d3-20_s4b}
\includegraphics[width=.49\linewidth]{Eff-sigstar-normal_n2d_d3-20_s4b}
\end{center}
\caption{\small  Left: $\ms^*(d,n,s)$ ({\color{Dandelion} $\blacklozenge$}),
$a^*(d,n,s)$ ({\color{Cerulean} $\bullet$}),
$\widetilde{a^*}(d,n,s)$ ({\color{green} $\blacktriangledown$}),
$M_{\psi,1}\widetilde{a^*}(d,n,s)$ ({\color{gray} {\tiny $\blacksquare$}}),
$\widehat{a^*_\ml}=\sqrt{3}/2$ ({\color{Orchid} $\blacktriangle$}) and
$\widehat{a^*}(d,n,s)$ ({\color{red} $\bigstar$}) as functions of $d$.
Right: efficiencies $D_{\mu,s}^{1/s}(\PP_a^{[n]})/D_{\mu,s}^{1/s}(\mu_{\ms^*}^{[n]})$ as functions of $d$ for the choices of $a$ indicated on the left panel ($n=2^d$, $s=4$).}
\label{F:normal_n2d_d3-20_s4}
\end{figure}
% tst_quantization_random_normal.m

\begin{remark}
From Corollary~\ref{Coro:astar-asymptotic} and Proposition~\ref{Prop:extreme-value-general-mu},all quantities  $a^*(d,n,s)$, $\widetilde{a^*}(d,n,s)$,
$M_{\psi,1}\widetilde{a^*}(d,n,s)$ and
%$M_{\psi,1}\widehat{a^*_\ml}$
$\widehat{a^*}(d,n,s)$  tend to
$\widehat{a^*_\ml}=\sqrt{3}/2$ when $d\to\infty$. Our numerical results suggest that, in the considered asymptotic regime, $\ms^*(d,n,s)$ also converges to the limiting value $\widehat{a^*_\ml} = \sqrt{3}/2$. However, a rigorous theoretical analysis is still required to confirm this observation. The same is true for the asymptotic behaviour of $b^*(d,n,s)$ in the optimised quantiser $\PP_{0,b^*}^{[n]}$ when $\mu$ is uniform in $\SB_d(1)$. From \eqref{eq:quantile2} in Proposition~\ref{th:dist}, we have to analyse the behaviour of $d(\ub,\Rb_n)=\min_{i=1,\ldots,n} \left[(\|\ub\|-R_i)^2 +4\,\|\ub\|\,R_i \, \zeta_i\right]$ when $n\to\infty$, where the random variables $R_i$ are $\zeta_i$ are i.i.d.\ and mutually independent, with $\zeta_i \eqd \beta_{\delta,\delta}$. The assumption that $R=\|X\|$ with the probability measure $\PP$ of $X$ satisfying a norm-concentration property like \eqref{norm-concentration} does not directly imply an extreme-value property for $d(\ub,\Rb_n)$. The cases where $\PP = \PP_a$ (uniform distribution on $\SB_d(a)$) and $\PP = \mu_\ms$ (multivariate normal distribution) already pose significant challenges.
\fin
\end{remark}


%___________________________________________________________________________
\section{Conclusions}~\label{S:conclusion}

We have demonstrated that, for a spherically symmetric measure $\mu$ in $\RR^d$, unless the sample size $n$ is extremely large, a random quantiser uniformly distributed on a sphere of suitable radius can significantly outperform a random quantiser whose distribution follows the asymptotic limit predicted by Zador’s theorem. The optimal radius can be determined numerically by minimising a triple integral, which can be evaluated with arbitrary precision. Additionally, an approximation of the optimal radius is available, derived from extreme-value theory. Since the variability of the distortion for such random quantisers rapidly diminishes as $n$ increases, this approach is highly practical for constructing high-performance quantisers.

While the restriction to spherically symmetric distributions may appear limiting, it provides a robust foundation for broader applications. For instance, quantising the uniform measure on the $d$-dimensional hypercube $[-1,1]^d$ presents a compelling and challenging problem, particularly in space-filling design for computer experiments (see, e.g., \cite{PronMul2012}).  While greedy quantisation offers an attractive approach for the incremental construction of designs in low dimensions \cite{LuschgyP2015,NogalesPR2021}, it becomes computationally infeasible for large $d$. In such cases, recent results from \cite{noonanhigh2024} allow the uniform measure to be approximated by a spherically symmetric distribution. Similarly, one can consider random quantisers distributed according to a product measure, which can itself be approximated by a spherically symmetric distribution. Preliminary findings then suggest that quantisers distributed on the vertices of a smaller hypercube exhibit promising performance, paving the way for further advances in this direction.

%___________________________________________________________________________
\bibliographystyle{plain}
\bibliography{high-d} %,springer_luc}



\end{document}




